\documentclass{article}
\usepackage{amsmath, amssymb, geometry}
\geometry{a4paper, margin=1in}

\title{Homework Chapter 1 Sets}
\author{Tashfeen Omran}
\date{September 2026}

\begin{document}

\maketitle

\section{Comprehensive Introduction, Context, and Prerequisites}

\subsection{Core Concepts: What is a Set?}
At its heart, mathematics is built upon the fundamental concept of a \textbf{set}. A set is simply a well-defined collection of distinct objects. These objects are called \textbf{elements} or \textbf{members} of the set. The term "well-defined" means that for any given object, it must be clear whether or not it belongs to the set.

Sets are typically denoted by capital letters (e.g., $A, B, S$), and their elements are enclosed within curly braces $\{ \}$. For example, the set $S = \{1, 2, 3\}$ consists of the numbers 1, 2, and 3. The order in which elements are listed does not matter, and elements are not repeated. So, $\{1, 2, 3\}$ is the same set as $\{3, 2, 1\}$ and $\{1, 1, 2, 3\}$.

\subsubsection{Notation}
\begin{itemize}
    \item $x \in A$: Denotes that $x$ is an element of set $A$.
    \item $x \notin A$: Denotes that $x$ is not an element of set $A$.
\end{itemize}

\subsubsection{Describing Sets}
\begin{itemize}
    \item \textbf{Listing Method:} Explicitly listing all elements (if the set is finite and small enough), e.g., $A = \{a, b, c\}$.
    \item \textbf{Set-Builder Notation:} Describing the elements by a property they satisfy, e.g., $B = \{x \mid x \text{ is an even integer}\}$. This reads "the set of all $x$ such that $x$ is an even integer." Sometimes a colon $:$ is used instead of a vertical bar $\mid$.
\end{itemize}

\subsubsection{Types of Sets}
\begin{itemize}
    \item \textbf{Finite Set:} A set with a countable number of elements, including the empty set.
    \item \textbf{Infinite Set:} A set with an uncountable number of elements, e.g., the set of real numbers.
    \item \textbf{Empty Set ($\emptyset$ or \{\}):} A unique set containing no elements. Its cardinality is 0.
\end{itemize}

\subsubsection{Cardinality}
The \textbf{cardinality} of a finite set $A$, denoted $|A|$, is the number of distinct elements in $A$. For example, if $A = \{1, 2, 3, 4\}$, then $|A| = 4$.

\subsubsection{Special Sets of Numbers}
\begin{itemize}
    \item $\mathbb{N}$: The set of \textbf{natural numbers} (positive integers), $\mathbb{N} = \{1, 2, 3, \dots\}$.
    \item $\mathbb{Z}$: The set of \textbf{integers} (positive, negative, and zero), $\mathbb{Z} = \{\dots, -2, -1, 0, 1, 2, \dots\}$.
    \item $\mathbb{Q}$: The set of \textbf{rational numbers}, $\mathbb{Q} = \{\frac{m}{n} \mid m, n \in \mathbb{Z}, n \neq 0\}$.
    \item $\mathbb{I}$: The set of \textbf{irrational numbers} (real numbers that are not rational), e.g., $\sqrt{2}, \pi, e$.
    \item $\mathbb{R}$: The set of \textbf{real numbers} ($\mathbb{Q} \cup \mathbb{I}$).
    \item $\mathbb{C}$: The set of \textbf{complex numbers}, $\mathbb{C} = \{a + bi \mid a, b \in \mathbb{R}, i = \sqrt{-1}\}$.
\end{itemize}

\subsubsection{Subsets and Power Sets}
\begin{itemize}
    \item \textbf{Subset ($A \subseteq B$):} Set $A$ is a subset of set $B$ if every element of $A$ is also an element of $B$.
    \item \textbf{Proper Subset ($A \subset B$):} Set $A$ is a proper subset of $B$ if $A \subseteq B$ and $A \neq B$.
    \item \textbf{Universal Set ($U$):} In any particular context, all sets under consideration are subsets of a larger set called the universal set.
    \item \textbf{Power Set ($P(A)$):} The set of all possible subsets of $A$. If $|A| = n$, then $|P(A)| = 2^n$.
\end{itemize}

\subsubsection{Set Operations}
\begin{itemize}
    \item \textbf{Union ($A \cup B$):} The set of all elements that are in $A$ or in $B$ (or both). $A \cup B = \{x \mid x \in A \text{ or } x \in B\}$.
    \item \textbf{Intersection ($A \cap B$):} The set of all elements that are in both $A$ and $B$. $A \cap B = \{x \mid x \in A \text{ and } x \in B\}$.
    \item \textbf{Disjoint Sets:} Two sets $A$ and $B$ are disjoint if $A \cap B = \emptyset$.
    \item \textbf{Set Difference ($A - B$ or $A \setminus B$):} The set of all elements that are in $A$ but not in $B$. $A - B = \{x \mid x \in A \text{ and } x \notin B\}$.
    \item \textbf{Complement ($\bar{A}$ or $A^c$):} The set of all elements in the universal set $U$ that are not in $A$. $\bar{A} = U - A = \{x \mid x \in U \text{ and } x \notin A\}$.
    \item \textbf{Cartesian Product ($A \times B$):} The set of all possible ordered pairs $(a, b)$ where $a \in A$ and $b \in B$. $A \times B = \{(a, b) \mid a \in A \text{ and } b \in B\}$. If $|A|=m$ and $|B|=n$, then $|A \times B|=m \cdot n$.
\end{itemize}

\subsubsection{Partitions}
A \textbf{partition} of a nonempty set $A$ is a collection of nonempty subsets of $A$ (called blocks or cells) such that:
\begin{enumerate}
    \item Every element of $A$ belongs to exactly one subset in the collection.
    \item The subsets are pairwise disjoint (their intersection is empty).
    \item The union of all subsets in the collection equals $A$.
\end{enumerate}

\subsection{Intuition and Derivation}
The intuition behind sets is simple: they are containers for objects. Just as we collect physical objects, sets allow us to collect mathematical objects with precision. There isn't a "derivation" of sets in the same way there is for a formula like arc length. Instead, set theory provides the foundational language and axioms upon which most of modern mathematics is built. The definitions of union, intersection, and other operations are direct translations of logical "and," "or," and "not" into operations on collections. For instance, the union $A \cup B$ is intuitive: if something is in $A$ or in $B$, it's in the combined collection. The power set's cardinality of $2^n$ stems from the combinatorial choice for each of $n$ elements: either it's in a given subset or it's not, leading to $2 \times 2 \times \dots \times 2$ ($n$ times) possibilities.

\subsection{Historical Context and Motivation}
The formal study of sets began in the late 19th century, primarily with the work of Georg Cantor. Before Cantor, mathematicians often used the idea of collections of objects informally. However, as mathematics became more rigorous, the need for a precise and foundational understanding of these collections became apparent. Cantor initially developed set theory to understand and classify different types of infinite numbers (cardinalities), which led to profound and sometimes counter-intuitive results (like there being "more" real numbers than natural numbers).

The initial "naive" set theory, while powerful, also led to paradoxes (most famously, Russell's Paradox). This motivated the development of axiomatic set theory (like Zermelo-Fraenkel set theory with the Axiom of Choice, ZFC), which provides a more robust and consistent foundation for mathematics, avoiding such contradictions. Set theory now serves as the universal language for mathematics, allowing for clear and unambiguous definitions of nearly all mathematical objects and structures.

\subsection{Key Formulas}
\begin{itemize}
    \item \textbf{Cardinality of a finite set $A$:} $|A|$
    \item \textbf{Cardinality of the Power Set of $A$:} $|P(A)| = 2^{|A|}$
    \item \textbf{Union:} $A \cup B = \{x \mid x \in A \text{ or } x \in B\}$
    \item \textbf{Intersection:} $A \cap B = \{x \mid x \in A \text{ and } x \in B\}$
    \item \textbf{Set Difference:} $A - B = \{x \mid x \in A \text{ and } x \notin B\}$
    \item \textbf{Complement (given universal set $U$):} $\bar{A} = U - A = \{x \mid x \in U \text{ and } x \notin A\}$
    \item \textbf{Cartesian Product:} $A \times B = \{(a, b) \mid a \in A \text{ and } b \in B\}$
    \item \textbf{Cardinality of Cartesian Product:} $|A \times B| = |A| \cdot |B|$
\end{itemize}

\subsection{Prerequisites}
To effectively understand and work with sets, the following foundational skills are essential:
\begin{itemize}
    \item \textbf{Basic Logic:} Understanding of logical connectives such as "and" ($\land$), "or" ($\lor$), "not" ($\neg$), and implications. These directly translate to set operations like intersection, union, and complement.
    \item \textbf{Number Systems:} Familiarity with different categories of numbers: natural numbers ($\mathbb{N}$), integers ($\mathbb{Z}$), rational numbers ($\mathbb{Q}$), real numbers ($\mathbb{R}$), and complex numbers ($\mathbb{C}$).
    \item \textbf{Algebraic Manipulation:} Solving linear and quadratic equations and inequalities to determine the elements of sets defined by conditions.
    \item \textbf{Interval Notation:} Understanding and manipulating intervals on the real number line (e.g., $(a, b)$, $[a, b]$, $(-\infty, a)$).
    \item \textbf{Absolute Value:} Solving equations and inequalities involving absolute values, which are frequently used in set definitions.
\end{itemize}

\section{Detailed Homework Solutions}

\subsection*{Chapter 1 Supplemental Exercises}

\subsubsection*{Problem 1.73}
The set $T = \{2k + 1 : k \in \mathbb{Z}\}$ can be described as $T = \{\dots, -3, -1, 1, 3, \dots\}$. Describe the following sets in a similar manner.

\paragraph*{Solution to 1.73 (a)}
$A = \{4k + 3 : k \in \mathbb{Z}\}$
Let's substitute a few integer values for $k$:
If $k = -2$, $4(-2) + 3 = -8 + 3 = -5$.
If $k = -1$, $4(-1) + 3 = -4 + 3 = -1$.
If $k = 0$, $4(0) + 3 = 0 + 3 = 3$.
If $k = 1$, $4(1) + 3 = 4 + 3 = 7$.
If $k = 2$, $4(2) + 3 = 8 + 3 = 11$.
So, $A = \{\dots, -5, -1, 3, 7, 11, \dots\}$.

\paragraph*{Solution to 1.73 (b)}
$B = \{5k - 1 : k \in \mathbb{Z}\}$
Let's substitute a few integer values for $k$:
If $k = -2$, $5(-2) - 1 = -10 - 1 = -11$.
If $k = -1$, $5(-1) - 1 = -5 - 1 = -6$.
If $k = 0$, $5(0) - 1 = 0 - 1 = -1$.
If $k = 1$, $5(1) - 1 = 5 - 1 = 4$.
If $k = 2$, $5(2) - 1 = 10 - 1 = 9$.
So, $B = \{\dots, -11, -6, -1, 4, 9, \dots\}$.

\subsubsection*{Problem 1.74}
Let $S = \{-10, -9, \dots, 9, 10\}$. Describe each of the following sets as $\{x \in S : p(x)\}$, where $p(x)$ is some condition on $x$.

\paragraph*{Solution to 1.74 (a)}
$A = \{-10, -9, \dots, -1, 1, \dots, 9, 10\}$
This set contains all elements of $S$ except $0$.
$A = \{x \in S : x \neq 0\}$.

\paragraph*{Solution to 1.74 (b)}
$B = \{-10, -9, \dots, -1, 0\}$
This set contains all elements of $S$ that are less than or equal to $0$.
$B = \{x \in S : x \le 0\}$.

\paragraph*{Solution to 1.74 (c)}
$C = \{-5, -4, \dots, 0, 1, \dots, 7\}$
This set contains elements of $S$ that are greater than or equal to $-5$ and less than or equal to $7$.
$C = \{x \in S : -5 \le x \le 7\}$.

\paragraph*{Solution to 1.74 (d)}
$D = \{-10, -9, \dots, 4, 6, 7, \dots, 10\}$
This set contains all elements of $S$ except $5$.
$D = \{x \in S : x \neq 5\}$.

\subsubsection*{Problem 1.75}
Describe each of the following sets by listing its elements within braces.

\paragraph*{Solution to 1.75 (a)}
$\{x \in \mathbb{Z} : x^3 - 4x = 0\}$
Factor the equation: $x(x^2 - 4) = 0 \implies x(x-2)(x+2) = 0$.
The integer solutions are $x = 0, x = 2, x = -2$.
So, the set is $\{-2, 0, 2\}$.

\paragraph*{Solution to 1.75 (b)}
$\{x \in \mathbb{R} : |x| = -1\}$
The absolute value of a real number is always non-negative ($|x| \ge 0$). There is no real number whose absolute value is $-1$.
So, the set is $\emptyset$.

\paragraph*{Solution to 1.75 (c)}
$\{m \in \mathbb{N} : 2 < m \le 5\}$
The natural numbers greater than 2 and less than or equal to 5 are $3, 4, 5$.
So, the set is $\{3, 4, 5\}$.

\paragraph*{Solution to 1.75 (d)}
$\{n \in \mathbb{N} : 0 \le n \le 3\}$
The natural numbers greater than or equal to 0 and less than or equal to 3. Since $\mathbb{N}$ usually refers to positive integers ($1, 2, 3, \dots$), $0 \notin \mathbb{N}$. If $\mathbb{N}$ includes $0$, it would be $\{0, 1, 2, 3\}$. Given the book's convention (page 16, $\mathbb{N}=\{1,2,3,\dots\}$), $0$ is not a natural number.
So, the set is $\{1, 2, 3\}$.

\paragraph*{Solution to 1.75 (e)}
$\{k \in \mathbb{Q} : k^2 - 4 = 0\}$
Factor the equation: $(k-2)(k+2) = 0$.
The solutions are $k = 2, k = -2$. Both are rational numbers.
So, the set is $\{-2, 2\}$.

\paragraph*{Solution to 1.75 (f)}
$\{k \in \mathbb{Z} : 9k^2 - 3 = 0\}$
$9k^2 = 3 \implies k^2 = \frac{3}{9} \implies k^2 = \frac{1}{3}$.
The solutions are $k = \sqrt{\frac{1}{3}}$ and $k = -\sqrt{\frac{1}{3}}$. These are irrational numbers and therefore not integers.
So, the set is $\emptyset$.

\paragraph*{Solution to 1.75 (g)}
$\{k \in \mathbb{Z} : 1 \le k^2 < 10\}$
We need integers $k$ whose square is between 1 (inclusive) and 10 (exclusive).
Possible squares are $1, 4, 9$.
If $k^2 = 1$, then $k = 1$ or $k = -1$.
If $k^2 = 4$, then $k = 2$ or $k = -2$.
If $k^2 = 9$, then $k = 3$ or $k = -3$.
So, the set is $\{-3, -2, -1, 1, 2, 3\}$.

\subsubsection*{Problem 1.76}
Determine the cardinality of each of the following sets.

\paragraph*{Solution to 1.76 (a)}
$A = \{1, 2, 3, \{1, 2, 3\}, 4, \{4\}\}$
The distinct elements are $1, 2, 3, \{1, 2, 3\}, 4, \{4\}$.
There are 6 distinct elements.
So, $|A| = 6$.

\paragraph*{Solution to 1.76 (b)}
$B = \{x \in \mathbb{R} : |x| = -1\}$
As determined in Problem 1.75 (b), this set is $\emptyset$.
So, $|B| = 0$.

\paragraph*{Solution to 1.76 (c)}
$C = \{m \in \mathbb{N} : 2 < m < 5\}$
The natural numbers greater than 2 and less than 5 are $3, 4$.
So, $|C| = 2$.

\paragraph*{Solution to 1.76 (d)}
$D = \{n \in \mathbb{N} : n < 0\}$
The natural numbers are $\{1, 2, 3, \dots\}$. There are no natural numbers less than 0.
So, the set is $\emptyset$.
$|D| = 0$.

\paragraph*{Solution to 1.76 (e)}
$E = \{k \in \mathbb{N} : 1 \le k^2 < 100\}$
We need natural numbers $k$ whose square is between 1 (inclusive) and 100 (exclusive).
$k^2 \ge 1 \implies k \ge 1$ (since $k \in \mathbb{N}$).
$k^2 < 100 \implies k < 10$ (since $k \in \mathbb{N}$).
So, $k \in \{1, 2, 3, 4, 5, 6, 7, 8, 9\}$.
$|E| = 9$.

\paragraph*{Solution to 1.76 (f)}
$F = \{k \in \mathbb{Z} : 1 \le k^2 < 100\}$
We need integers $k$ whose square is between 1 (inclusive) and 100 (exclusive).
The possible values for $k^2$ are $1, 4, 9, \dots, 81$.
This means $k$ can be $\pm 1, \pm 2, \dots, \pm 9$.
So, $k \in \{-9, -8, \dots, -1, 1, \dots, 8, 9\}$.
$|F| = 9 \times 2 = 18$.

\subsubsection*{Problem 1.77}
For $A = \{-1, 0, 1\}$ and $B = \{x, y\}$, determine $A \times B$.

\paragraph*{Solution to 1.77}
$A \times B = \{(-1, x), (-1, y), (0, x), (0, y), (1, x), (1, y)\}$.

\subsubsection*{Problem 1.78}
Let $U = \{1, 2, 3\}$ be the universal set and let $A = \{1, 2\}$, $B = \{2, 3\}$ and $C = \{1, 3\}$. Determine the following.

\paragraph*{Solution to 1.78 (a)}
$(A \cup B) - (B \cap C)$
First, $A \cup B = \{1, 2\} \cup \{2, 3\} = \{1, 2, 3\}$.
Next, $B \cap C = \{2, 3\} \cap \{1, 3\} = \{3\}$.
Then, $(A \cup B) - (B \cap C) = \{1, 2, 3\} - \{3\} = \{1, 2\}$.

\paragraph*{Solution to 1.78 (b)}
$\bar{A}$
The universal set is $U = \{1, 2, 3\}$ and $A = \{1, 2\}$.
$\bar{A} = U - A = \{1, 2, 3\} - \{1, 2\} = \{3\}$.

\paragraph*{Solution to 1.78 (c)}
$B \cup \bar{C}$
First, $\bar{C} = U - C = \{1, 2, 3\} - \{1, 3\} = \{2\}$.
Then, $B \cup \bar{C} = \{2, 3\} \cup \{2\} = \{2, 3\}$.

\paragraph*{Solution to 1.78 (d)}
$A \times B$
$A = \{1, 2\}$ and $B = \{2, 3\}$.
$A \times B = \{(1, 2), (1, 3), (2, 2), (2, 3)\}$.

\subsubsection*{Problem 1.79}
Let $A = \{1, 2, \dots, 10\}$. Give an example of two sets $S$ and $B$ such that $S \subseteq P(A)$, $|S| = 4$, $B \in S$ and $|B| = 2$.

\paragraph*{Solution to 1.79}
We need $S$ to be a set of subsets of $A$. $B$ must be one of these subsets, and $B$ itself must have 2 elements.
Let's choose $B = \{1, 2\}$. Since $B \in S$, $S$ must contain $\{1, 2\}$.
We need $|S|=4$, so $S$ must contain three other subsets of $A$. Let's pick some arbitrary subsets of $A$.
Let $B = \{1, 2\}$.
Let $S = \{\{1, 2\}, \{3, 4\}, \{5, 6\}, \{7\}\}$.
Here, $S \subseteq P(A)$ (since all elements of $S$ are subsets of $A$).
$|S| = 4$.
$B = \{1, 2\} \in S$.
$|B| = |\{1, 2\}| = 2$.
This example satisfies all conditions.

\subsubsection*{Problem 1.80}
For $A = \{1\}$ and $C = \{1, 2\}$, give an example of a set $B$ such that $P(A) \subset B \subset P(C)$.

\paragraph*{Solution to 1.80}
First, determine $P(A)$ and $P(C)$.
$A = \{1\}$. Subsets of $A$ are $\emptyset$ and $\{1\}$. So $P(A) = \{\emptyset, \{1\}\}$.
$C = \{1, 2\}$. Subsets of $C$ are $\emptyset, \{1\}, \{2\}, \{1, 2\}$. So $P(C) = \{\emptyset, \{1\}, \{2\}, \{1, 2\}\}$.
We need a set $B$ such that $P(A) \subset B \subset P(C)$. This means $B$ must contain all elements of $P(A)$, must be a proper subset of $P(C)$, and must contain at least one element not in $P(A)$.
We can take $B = P(A) \cup \{\text{some element from } P(C) \text{ not in } P(A)\}$.
Elements of $P(C)$ not in $P(A)$ are $\{2\}$ and $\{1, 2\}$.
Let's choose $B = \{\emptyset, \{1\}, \{2\}\}$.
Here, $P(A) = \{\emptyset, \{1\}\} \subset B = \{\emptyset, \{1\}, \{2\}\}$.
And $B = \{\emptyset, \{1\}, \{2\}\} \subset P(C) = \{\emptyset, \{1\}, \{2\}, \{1, 2\}\}$.
This satisfies the conditions. Another example could be $B = \{\emptyset, \{1\}, \{1, 2\}\}$.

\subsubsection*{Problem 1.81}
Give examples of two sets $A$ and $B$ such that $A \cap P(A) \in B$ and $P(A) \subset A \cup B$.

\paragraph*{Solution to 1.81}
Let's try a simple set for $A$.
Let $A = \{1\}$.
Then $P(A) = \{\emptyset, \{1\}\}$.
$A \cap P(A) = \{1\} \cap \{\emptyset, \{1\}\} = \{1\}$.
So, we need $\{1\} \in B$.
Now, consider $P(A) \subset A \cup B$. This means $\{\emptyset, \{1\}\} \subset \{1\} \cup B$.
This implies that $\emptyset$ must be an element of $A \cup B$. Since $\emptyset \notin A$, it must be that $\emptyset \in B$.
So, we have two conditions for $B$: $\{1\} \in B$ and $\emptyset \in B$.
Let's choose a $B$ that satisfies these conditions, for instance, $B = \{\emptyset, \{1\}\}$.
Let's check the conditions:
1. $A \cap P(A) = \{1\}$. Is $\{1\} \in B$? Yes, $\{1\} \in \{\emptyset, \{1\}\}$.
2. $P(A) = \{\emptyset, \{1\}\}$. Is $P(A) \subset A \cup B$?
   $A \cup B = \{1\} \cup \{\emptyset, \{1\}\} = \{\emptyset, 1, \{1\}\}$.
   Is $\{\emptyset, \{1\}\} \subset \{\emptyset, 1, \{1\}\}$? Yes, because every element of $P(A)$ ($\emptyset$ and $\{1\}$) is an element of $A \cup B$, and $A \cup B$ has an additional element ($1$).
So, $A = \{1\}$ and $B = \{\emptyset, \{1\}\}$ is a valid example.

\subsubsection*{Problem 1.82}
Which of the following sets are equal?
$A = \{n \in \mathbb{Z} : -4 \le n \le 4\}$
$B = \{x \in \mathbb{N} : 2x + 2 = 0\}$
$C = \{x \in \mathbb{Z} : 3x - 2 = 0\}$
$D = \{x \in \mathbb{Z} : x^3 = 4x\}$
$E = \{-2, 0, 2\}$

\paragraph*{Solution to 1.82}
Let's list the elements of each set:
\begin{itemize}
    \item $A = \{n \in \mathbb{Z} : -4 \le n \le 4\}$
    $A = \{-4, -3, -2, -1, 0, 1, 2, 3, 4\}$.
    \item $B = \{x \in \mathbb{N} : 2x + 2 = 0\}$
    $2x + 2 = 0 \implies 2x = -2 \implies x = -1$.
    Since $\mathbb{N} = \{1, 2, 3, \dots\}$, $-1 \notin \mathbb{N}$.
    $B = \emptyset$.
    \item $C = \{x \in \mathbb{Z} : 3x - 2 = 0\}$
    $3x - 2 = 0 \implies 3x = 2 \implies x = \frac{2}{3}$.
    Since $\frac{2}{3}$ is not an integer, $x \notin \mathbb{Z}$.
    $C = \emptyset$.
    \item $D = \{x \in \mathbb{Z} : x^3 = 4x\}$
    $x^3 - 4x = 0 \implies x(x^2 - 4) = 0 \implies x(x-2)(x+2) = 0$.
    The integer solutions are $x = 0, x = 2, x = -2$.
    $D = \{-2, 0, 2\}$.
    \item $E = \{-2, 0, 2\}$.
\end{itemize}
Comparing the sets, we find that $D = E$. Also, $B = C$.
So, $D$ and $E$ are equal. $B$ and $C$ are equal.

\subsubsection*{Problem 1.83}
Let $A$ and $B$ be subsets of some unknown universal set $U$. Suppose that $\bar{A} = \{3, 8, 9\}$, $\bar{A} \cup B = \{1, 2\}$, $B - A = \{8\}$ and $A \cap B = \{5, 7\}$. Determine $U, A$ and $B$.

\paragraph*{Solution to 1.83}
We are given:
1. $\bar{A} = \{3, 8, 9\}$
2. $\bar{A} \cup B = \{1, 2\}$
3. $B - A = \{8\}$
4. $A \cap B = \{5, 7\}$

From (1), $A$ is the complement of $\bar{A}$ with respect to $U$. The elements $3, 8, 9$ are in $\bar{A}$, so they are NOT in $A$.
From (4), $A \cap B = \{5, 7\}$. This means $5 \in A$, $7 \in A$, $5 \in B$, $7 \in B$.

Now consider (3), $B - A = \{8\}$. This means $8 \in B$ and $8 \notin A$.

Let's use (2), $\bar{A} \cup B = \{1, 2\}$.
We know $\bar{A} = \{3, 8, 9\}$.
So, $\{3, 8, 9\} \cup B = \{1, 2\}$.
For this equality to hold, any element in $\{3, 8, 9\}$ must also be in $\{1, 2\}$ if it's not removed by some property of $B$, which is impossible because $B$ adds elements. This is a contradiction, as $3 \in \bar{A} \cup B$ but $3 \notin \{1, 2\}$. Let me re-read the problem statement for $\bar{A} \cup B = \{1, 2\}$.
Ah, $\bar{A} \cup B = \{1, 2\}$ means that all elements $x$ such that $x \in \bar{A}$ or $x \in B$ must be $1$ or $2$.
However, we know $3 \in \bar{A}$, $8 \in \bar{A}$, $9 \in \bar{A}$. This implies $3 \in \bar{A} \cup B$, $8 \in \bar{A} \cup B$, $9 \in \bar{A} \cup B$.
But $\bar{A} \cup B = \{1, 2\}$ does not contain $3, 8, 9$.
This indicates an inconsistency in the problem statement as written. Let me assume the OCR might have a slight error, or interpret this as a constraint *on* $B$.
If $\bar{A} \cup B = \{1, 2\}$, then the elements $3, 8, 9$ cannot be in $\bar{A}$. This directly contradicts $\bar{A} = \{3, 8, 9\}$.

Let's assume the problem meant $A \cup \bar{B} = \{1, 2\}$ or similar, or that $U$ is limited.
Let's restart by carefully compiling all information we have about elements:
\begin{itemize}
    \item From (1) $\bar{A} = \{3, 8, 9\}$:
        *   Elements definitely in $U$: $3, 8, 9$.
        *   Elements definitely $\notin A$: $3, 8, 9$.
    \item From (4) $A \cap B = \{5, 7\}$:
        *   Elements definitely in $A$: $5, 7$.
        *   Elements definitely in $B$: $5, 7$.
    \item From (3) $B - A = \{8\}$:
        *   Elements definitely in $B$: $8$.
        *   Elements definitely $\notin A$: $8$. (This is consistent with $\bar{A} = \{3, 8, 9\}$).
\end{itemize}
So far:
$A$ contains $\{5, 7\}$. $A$ does not contain $3, 8, 9$.
$B$ contains $\{5, 7, 8\}$. $B$ does not contain $3, 9$ (since $3, 9 \in \bar{A}$ and $A \cap B = \{5,7\}$ so elements in $B$ not in $A$ are $8$, therefore $3, 9 \notin B$).

Let's use this information for (2): $\bar{A} \cup B = \{1, 2\}$.
From our current understanding:
$\bar{A} = \{3, 8, 9\}$.
$B = \{5, 7, 8\}$ (assuming these are all the elements in $B$ from the given info).
Then $\bar{A} \cup B = \{3, 8, 9\} \cup \{5, 7, 8\} = \{3, 5, 7, 8, 9\}$.
But this should be equal to $\{1, 2\}$.
This is a direct contradiction. The problem statement as provided contains conflicting information.

Let's reconsider the possibility of a typo in the problem source itself, or if "A-bar U B = {1,2}" implies that the universal set must be limited to $\{1,2,3,5,7,8,9\}$ for those to be the only elements.
Let's assume there is no typo and try to derive $U$ first from the given sets of elements.
Elements explicitly mentioned: $3, 8, 9, 1, 2, 5, 7$.
So $U$ must contain at least these elements: $\{1, 2, 3, 5, 7, 8, 9\}$.

From $\bar{A} = \{3, 8, 9\}$, we know $A$ cannot contain $3, 8, 9$.
From $A \cap B = \{5, 7\}$, we know $A$ contains $5, 7$. So $A$ is at least $\{5, 7\}$.
From $B - A = \{8\}$, we know $B$ contains $8$, and $A$ does not contain $8$.

Now, let's use $\bar{A} \cup B = \{1, 2\}$.
Since $3 \in \bar{A}$, it must be that $3 \in \{1, 2\}$, which is false.
This confirms a definite contradiction in the problem statement. It's impossible for these conditions to hold simultaneously.

\textit{Given the clear contradiction, I must state that the problem as written has no consistent solution.}
Assuming a common pattern, one of the values of $\bar{A}$ should be $1$ or $2$. Or $A \cup \bar{B}$ was intended.
If I were forced to make a "best effort" interpretation, I'd assume $U$ only contains the numbers mentioned, and try to resolve.
Let $U = \{1, 2, 3, 5, 7, 8, 9\}$.
From $\bar{A} = \{3, 8, 9\}$, $A = U - \bar{A} = \{1, 2, 5, 7\}$.
From $A \cap B = \{5, 7\}$, this is consistent with $A$. So $B$ must contain $5, 7$.
From $B - A = \{8\}$, $B$ must contain $8$ and $B$ does not contain elements from $A$ other than $8$.
So, $B$ must contain $5, 7, 8$.
Now, check $\bar{A} \cup B = \{1, 2\}$.
$\bar{A} = \{3, 8, 9\}$.
$B = \{5, 7, 8\}$.
$\bar{A} \cup B = \{3, 8, 9\} \cup \{5, 7, 8\} = \{3, 5, 7, 8, 9\}$.
This is NOT $\{1, 2\}$.

Conclusion: The problem statement contains an inherent contradiction. No sets $U, A, B$ can satisfy all given conditions simultaneously.

\subsubsection*{Problem 1.84}
Let $I$ denote the interval $[0, \infty)$. For each $r \in I$, define
$A_r = \{(x, y) \in \mathbb{R} \times \mathbb{R} : x^2 + y^2 = r^2\}$
$B_r = \{(x, y) \in \mathbb{R} \times \mathbb{R} : x^2 + y^2 < r^2\}$
$C_r = \{(x, y) \in \mathbb{R} \times \mathbb{R} : x^2 + y^2 > r^2\}$.

\paragraph*{Solution to 1.84 (a)}
Determine $\bigcup_{r \in I} A_r$ and $\bigcap_{r \in I} A_r$.
$A_r$ represents a circle centered at the origin with radius $r$.
$\bigcup_{r \in I} A_r$: This is the union of all circles centered at the origin with radius $r \ge 0$.
The point $(0,0)$ is in $A_0$ (since $0^2+0^2=0^2$). Any other point $(x,y)$ such that $x^2+y^2=k^2$ for some $k>0$ belongs to $A_k$.
So, the union covers all points $(x,y)$ in the plane such that $x^2+y^2 \ge 0$. This is the entire $\mathbb{R} \times \mathbb{R}$ plane.
$\bigcup_{r \in I} A_r = \mathbb{R} \times \mathbb{R}$.

$\bigcap_{r \in I} A_r$: This is the intersection of all circles centered at the origin with radius $r \ge 0$.
For a point $(x,y)$ to be in the intersection, it must satisfy $x^2+y^2 = r^2$ for *every* $r \in I$.
This is only possible if $r^2$ is always the same value, which is not true as $r$ varies over $[0, \infty)$.
The only point common to all $A_r$ is if $x^2+y^2=0$, meaning $(x,y)=(0,0)$, which only satisfies $A_0$. For any $r>0$, $(0,0) \notin A_r$.
Thus, there is no point common to all circles except for $(0,0)$ which is only in $A_0$.
So, $\bigcap_{r \in I} A_r = \emptyset$.

\paragraph*{Solution to 1.84 (b)}
Determine $\bigcup_{r \in I} B_r$ and $\bigcap_{r \in I} B_r$.
$B_r$ represents the open disk centered at the origin with radius $r$ (excluding the boundary).
$\bigcup_{r \in I} B_r$: This is the union of all open disks centered at the origin with radius $r \ge 0$.
Note that $B_0 = \{(x,y) : x^2+y^2 < 0\} = \emptyset$.
For any point $(x,y) \in \mathbb{R} \times \mathbb{R}$ where $(x,y) \neq (0,0)$, let $k^2 = x^2+y^2$. We can choose $r > k$. Then $(x,y) \in B_r$.
So, the union covers all points in $\mathbb{R} \times \mathbb{R}$.
$\bigcup_{r \in I} B_r = \mathbb{R} \times \mathbb{R}$.

$\bigcap_{r \in I} B_r$: This is the intersection of all open disks centered at the origin with radius $r \ge 0$.
For a point $(x,y)$ to be in the intersection, it must satisfy $x^2+y^2 < r^2$ for *every* $r \in I$.
The only point that can satisfy this is $(0,0)$, because if $(x,y) \neq (0,0)$, let $k^2 = x^2+y^2$. Then for any $r$ such that $r^2 \le k^2$, $(x,y) \notin B_r$.
If $(x,y) = (0,0)$, then $0^2+0^2 < r^2$ for all $r>0$. But $B_0 = \emptyset$. This means the intersection does not include $(0,0)$ if we consider $r=0$.
However, the common convention for such intersections would generally include $(0,0)$ for $r>0$.
Let's consider the definition: $r \in I = [0, \infty)$.
For $r=0$, $B_0 = \emptyset$. The intersection of any collection of sets that includes $\emptyset$ is $\emptyset$.
So, $\bigcap_{r \in I} B_r = \emptyset$.

\paragraph*{Solution to 1.84 (c)}
Determine $\bigcup_{r \in I} C_r$ and $\bigcap_{r \in I} C_r$.
$C_r$ represents the set of points outside the open disk of radius $r$.
$\bigcup_{r \in I} C_r$: This is the union of all sets of points outside open disks of radius $r \ge 0$.
$C_0 = \{(x,y) : x^2+y^2 > 0\}$. This is $\mathbb{R} \times \mathbb{R} - \{(0,0)\}$.
Since $C_0$ is part of the union, the union will be at least $C_0$.
If $(x,y) = (0,0)$, it is not in any $C_r$ for $r \ge 0$ since $0^2+0^2 > r^2$ is never true.
So, $\bigcup_{r \in I} C_r = \mathbb{R} \times \mathbb{R} - \{(0,0)\}$.

$\bigcap_{r \in I} C_r$: This is the intersection of all sets of points outside open disks of radius $r \ge 0$.
For a point $(x,y)$ to be in the intersection, it must satisfy $x^2+y^2 > r^2$ for *every* $r \in I$.
This is impossible. If $(x,y)$ is any fixed point, let $k^2 = x^2+y^2$. We can always choose an $r \in I$ such that $r^2 > k^2$ (e.g., $r = k+1$). Then $(x,y) \notin C_r$.
Thus, no point can satisfy the condition for all $r \in I$.
So, $\bigcap_{r \in I} C_r = \emptyset$.

\subsubsection*{Problem 1.85}
Give an example of four sets $A_1, A_2, A_3, A_4$ such that $|A_i \cap A_j| = |i - j|$ for every two integers $i$ and $j$ with $1 \le i \le j \le 4$.

\paragraph*{Solution to 1.85}
This problem asks us to construct sets with specific intersection cardinalities.
Let's list the required cardinalities for $|A_i \cap A_j|$:
\begin{itemize}
    \item $|A_1 \cap A_1| = |1-1| = 0$. This is unusual. Usually, $|A_i \cap A_i| = |A_i|$. This means $A_i$ must be $\emptyset$. But this seems to contradict other conditions.
    The phrasing "$|A_i \cap A_j| = |i-j|$" usually refers to $i \ne j$.
    If $|A_i \cap A_i| = 0$, then $A_i = \emptyset$. If $A_1 = \emptyset$, then $|A_1 \cap A_2| = |\emptyset \cap A_2| = 0$. But $|1-2| = 1$.
    So, this problem must imply $i \neq j$. If $i=j$, the condition $|A_i \cap A_i| = |i-i|=0$ implies $A_i = \emptyset$, which is not usually the intent for such problems. Let's assume the condition for $i=j$ is $|A_i \cap A_i| = |A_i|$ and the given condition applies for $i \ne j$. However, the wording "for every two integers $i$ and $j$ with $1 \le i \le j \le 4$" includes $i=j$. This is a slight ambiguity.

    Let's strictly follow the wording and assume $A_i$ must be empty based on $|A_i \cap A_i| = 0$.
    If $A_1 = \emptyset$, then $|A_1 \cap A_2| = 0$. But $|1-2|=1$. So $0=1$, which is impossible.
    Therefore, the interpretation of the problem statement for $i=j$ must be different or there is an implicit understanding that $i \ne j$.
    Let's assume the intent is for $i \ne j$. Then, the constraints are:
    \begin{itemize}
        \item $|A_1 \cap A_2| = |1-2| = 1$
        \item $|A_1 \cap A_3| = |1-3| = 2$
        \item $|A_1 \cap A_4| = |1-4| = 3$
        \item $|A_2 \cap A_3| = |2-3| = 1$
        \item $|A_2 \cap A_4| = |2-4| = 2$
        \item $|A_3 \cap A_4| = |3-4| = 1$
    \end{itemize}
    We also need to implicitly ensure that $|A_i|$ itself is large enough to contain all these intersection elements.

Let's use a Venn diagram approach for this. Assign unique elements for each intersection.
Let $A_1 = \{a, b, c, d, e, f\}$
Let $A_2 = \{a, g, h, i\}$
Let $A_3 = \{b, g, j, k, l\}$
Let $A_4 = \{c, h, j, m, n, o\}$

Let's define elements such that they satisfy the intersection cardinalities.
Let $S_{ij}$ be the set of elements unique to the intersection $A_i \cap A_j$.
$|A_1 \cap A_2| = 1$. Let $x_{12}$ be this element.
$|A_1 \cap A_3| = 2$. Let $x_{13a}, x_{13b}$ be these elements.
$|A_1 \cap A_4| = 3$. Let $x_{14a}, x_{14b}, x_{14c}$ be these elements.
$|A_2 \cap A_3| = 1$. Let $x_{23}$ be this element.
$|A_2 \cap A_4| = 2$. Let $x_{24a}, x_{24b}$ be these elements.
$|A_3 \cap A_4| = 1$. Let $x_{34}$ be this element.

Now we need to consider triple intersections. $A_1 \cap A_2 \cap A_3$, etc.
$A_1 \cap A_2 \cap A_3$: must contain elements that are also in $A_1 \cap A_2$, $A_1 \cap A_3$, $A_2 \cap A_3$.
Let's assign specific elements to each intersection.
$A_1 \cap A_2 = \{1\}$
$A_1 \cap A_3 = \{2, 3\}$
$A_1 \cap A_4 = \{4, 5, 6\}$
$A_2 \cap A_3 = \{7\}$
$A_2 \cap A_4 = \{8, 9\}$
$A_3 \cap A_4 = \{10\}$

Now check triple intersections:
$A_1 \cap A_2 \cap A_3$: Elements in $\{1\}, \{2,3\}, \{7\}$. No common elements. So $\emptyset$.
$A_1 \cap A_2 \cap A_4$: Elements in $\{1\}, \{4,5,6\}, \{8,9\}$. No common elements. So $\emptyset$.
$A_1 \cap A_3 \cap A_4$: Elements in $\{2,3\}, \{4,5,6\}, \{10\}$. No common elements. So $\emptyset$.
$A_2 \cap A_3 \cap A_4$: Elements in $\{7\}, \{8,9\}, \{10\}$. No common elements. So $\emptyset$.
This is a valid choice of elements for the pairwise intersections.

Now, define the sets $A_i$:
$A_1 = \{1, 2, 3, 4, 5, 6\}$ (contains all elements from its pairwise intersections)
$A_2 = \{1, 7, 8, 9\}$
$A_3 = \{2, 3, 7, 10\}$
$A_4 = \{4, 5, 6, 8, 9, 10\}$

Let's check the cardinalities:
$|A_1 \cap A_2| = |\{1\}| = 1$. (Correct)
$|A_1 \cap A_3| = |\{2, 3\}| = 2$. (Correct)
$|A_1 \cap A_4| = |\{4, 5, 6\}| = 3$. (Correct)
$|A_2 \cap A_3| = |\{7\}| = 1$. (Correct)
$|A_2 \cap A_4| = |\{8, 9\}| = 2$. (Correct)
$|A_3 \cap A_4| = |\{10\}| = 1$. (Correct)

This example satisfies all the conditions for $i \ne j$.

\subsubsection*{Problem 1.86}
(a) Give an example of two problems suggested by Exercise 1.85 (above).
(b) Solve one of the problems in (a).

\paragraph*{Solution to 1.86}
(a) Two problems suggested by Exercise 1.85:
\begin{enumerate}
    \item Give an example of three sets $A_1, A_2, A_3$ such that $|A_i \cap A_j| = \max(i, j) - \min(i, j) + 1$ for every two integers $i$ and $j$ with $1 \le i \le j \le 3$.
    \item Give an example of four sets $A_1, A_2, A_3, A_4$ such that $|A_i \cup A_j| = i+j$ for every two integers $i$ and $j$ with $1 \le i \le j \le 4$.
\end{enumerate}

(b) Solve problem 1: Give an example of three sets $A_1, A_2, A_3$ such that $|A_i \cap A_j| = \max(i, j) - \min(i, j) + 1$ for $1 \le i \le j \le 3$.
Let's list the required cardinalities:
\begin{itemize}
    \item $|A_1 \cap A_1| = \max(1,1) - \min(1,1) + 1 = 1-1+1 = 1$. So $|A_1|=1$.
    \item $|A_2 \cap A_2| = \max(2,2) - \min(2,2) + 1 = 2-2+1 = 1$. So $|A_2|=1$.
    \item $|A_3 \cap A_3| = \max(3,3) - \min(3,3) + 1 = 3-3+1 = 1$. So $|A_3|=1$.
    \item $|A_1 \cap A_2| = \max(1,2) - \min(1,2) + 1 = 2-1+1 = 2$. This contradicts $|A_1|=1$ and $|A_2|=1$. If $|A_1|=1$ and $|A_2|=1$, then $|A_1 \cap A_2|$ can be at most 1 (if $A_1=A_2$) or 0 (if $A_1 \ne A_2$). It cannot be 2.

This problem formulation, similar to the original 1.85, leads to a contradiction if interpreted for $i=j$.
Therefore, I must again assume the condition "$|A_i \cap A_j| = |i-j|$" (or the modified version here) implicitly refers to $i \ne j$.
Let's reformulate Problem 1 to remove this ambiguity.

\textbf{Revised Problem 1 (from 1.86 (a))}:
Give an example of three sets $A_1, A_2, A_3$ such that:
\begin{itemize}
    \item $|A_1 \cap A_2| = 1$
    \item $|A_1 \cap A_3| = 2$
    \item $|A_2 \cap A_3| = 1$
\end{itemize}
(This is simply problem 1.85 restricted to three sets).

\textbf{Solution to Revised Problem 1 (from 1.86 (b))}:
Let's assign elements.
$A_1 \cap A_2 = \{a\}$
$A_1 \cap A_3 = \{b, c\}$
$A_2 \cap A_3 = \{d\}$

For triple intersection, $A_1 \cap A_2 \cap A_3$.
Elements in $\{a\}, \{b, c\}, \{d\}$. No common elements. So $\emptyset$.

Now define the sets:
$A_1 = \{a, b, c\}$
$A_2 = \{a, d\}$
$A_3 = \{b, c, d\}$

Check cardinalities:
$|A_1 \cap A_2| = |\{a\}| = 1$. (Correct)
$|A_1 \cap A_3| = |\{b, c\}| = 2$. (Correct)
$|A_2 \cap A_3| = |\{d\}| = 1$. (Correct)
This example satisfies the conditions.

\subsubsection*{Problem 1.87}
Let $A = \{1, 2, 3\}$, $B = \{1, 2, 3, 4\}$ and $C = \{1, 2, 3, 4, 5\}$. For the sets $S$ and $T$ described below, explain whether $|S| < |T|$, $|S| > |T|$ or $|S| = |T|$.

\paragraph*{Solution to 1.87 (a)}
Let $B$ be the universal set and let $S$ be the set of all subsets $X$ of $B$ for which $|X| \neq |X|$. Let $T$ be the set of 2-element subsets of $C$.

First, let's analyze $S$. $S = \{X \subseteq B : |X| \neq |X|\}$.
The cardinality of any set is always equal to itself, i.e., $|X| = |X|$ is always true.
Therefore, the condition $|X| \neq |X|$ is never true.
So, $S$ is the empty set. $S = \emptyset$.
Thus, $|S| = 0$.

Next, let's analyze $T$. $T = \{\text{2-element subsets of } C\}$.
$C = \{1, 2, 3, 4, 5\}$. So $|C| = 5$.
The number of 2-element subsets of a set with 5 elements is given by the combination formula $\binom{n}{k} = \frac{n!}{k!(n-k)!}$.
Here, $n=5, k=2$.
$|T| = \binom{5}{2} = \frac{5!}{2!(5-2)!} = \frac{5!}{2!3!} = \frac{5 \times 4 \times 3 \times 2 \times 1}{(2 \times 1)(3 \times 2 \times 1)} = \frac{120}{12} = 10$.
So, $|T| = 10$.

Comparing $|S|$ and $|T|$: $0 < 10$.
Therefore, $|S| < |T|$.

\paragraph*{Solution to 1.87 (b)}
Let $S$ be the set of all partitions of the set $A$ and let $T$ be the set of 4-element subsets of $C$.

First, let's analyze $S$. $S = \{\text{partitions of } A\}$.
$A = \{1, 2, 3\}$. So $|A| = 3$.
The number of partitions of a set with $n$ elements is given by the Bell number $B_n$.
$B_3$:
Partitions of $\{1, 2, 3\}$ are:
\begin{itemize}
    \item $\{\{1, 2, 3\}\}$ (1 partition)
    \item $\{\{1, 2\}, \{3\}\}$
    \item $\{\{1, 3\}, \{2\}\}$
    \item $\{\{2, 3\}, \{1\}\}$ (3 partitions)
    \item $\{\{1\}, \{2\}, \{3\}\}$ (1 partition)
\end{itemize}
So, $|S| = B_3 = 1 + 3 + 1 = 5$.

Next, let's analyze $T$. $T = \{\text{4-element subsets of } C\}$.
$C = \{1, 2, 3, 4, 5\}$. So $|C| = 5$.
The number of 4-element subsets of a set with 5 elements is $\binom{5}{4}$.
$|T| = \binom{5}{4} = \frac{5!}{4!(5-4)!} = \frac{5!}{4!1!} = 5$.
So, $|T| = 5$.

Comparing $|S|$ and $|T|$: $5 = 5$.
Therefore, $|S| = |T|$.

\paragraph*{Solution to 1.87 (c)}
Let $S = \{(b, a) : b \in B, a \in A, a+b \text{ is odd}\}$ and let $T$ be the set of all nonempty proper subsets of $A$.

First, let's analyze $S$. $S = \{(b, a) : b \in B, a \in A, a+b \text{ is odd}\}$.
$A = \{1, 2, 3\}$.
$B = \{1, 2, 3, 4\}$.
For $a+b$ to be odd, one must be odd and the other must be even.
Odd elements in $A$: $\{1, 3\}$. Even elements in $A$: $\{2\}$.
Odd elements in $B$: $\{1, 3\}$. Even elements in $B$: $\{2, 4\}$.

Cases for $(b, a)$ where $a+b$ is odd:
\begin{itemize}
    \item $a$ is odd, $b$ is even:
        *   $a=1$: $(2,1), (4,1)$
        *   $a=3$: $(2,3), (4,3)$
    \item $a$ is even, $b$ is odd:
        *   $a=2$: $(1,2), (3,2)$
\end{itemize}
So, $S = \{(2,1), (4,1), (2,3), (4,3), (1,2), (3,2)\}$.
$|S| = 6$.

Next, let's analyze $T$. $T = \{\text{nonempty proper subsets of } A\}$.
$A = \{1, 2, 3\}$. So $|A| = 3$.
The total number of subsets of $A$ is $|P(A)| = 2^{|A|} = 2^3 = 8$.
These subsets are: $\emptyset, \{1\}, \{2\}, \{3\}, \{1,2\}, \{1,3\}, \{2,3\}, \{1,2,3\}$.
Nonempty subsets exclude $\emptyset$. So there are $8-1=7$ nonempty subsets.
Proper subsets exclude $A$ itself (i.e., $\{1,2,3\}$). So there are $7-1=6$ nonempty proper subsets.
$T = \{\{1\}, \{2\}, \{3\}, \{1,2\}, \{1,3\}, \{2,3\}\}$.
$|T| = 6$.

Comparing $|S|$ and $|T|$: $6 = 6$.
Therefore, $|S| = |T|$.

\subsubsection*{Problem 1.88}
Give an example of a set $A = \{1, 2, \dots, k\}$ for a smallest $k \in \mathbb{N}$ containing subsets $A_1, A_2, A_3$ such that $|A_i - A_j| = |A_j - A_i| = |i - j|$ for every two integers $i$ and $j$ with $1 \le i < j \le 3$.

\paragraph*{Solution to 1.88}
We need $A_1, A_2, A_3$ to be subsets of $A = \{1, 2, \dots, k\}$.
Conditions for $1 \le i < j \le 3$:
\begin{itemize}
    \item $i=1, j=2$: $|A_1 - A_2| = |1-2| = 1$ and $|A_2 - A_1| = |1-2| = 1$.
    \item $i=1, j=3$: $|A_1 - A_3| = |1-3| = 2$ and $|A_3 - A_1| = |1-3| = 2$.
    \item $i=2, j=3$: $|A_2 - A_3| = |2-3| = 1$ and $|A_3 - A_2| = |2-3| = 1$.
\end{itemize}
Let's represent the sets using unique elements for each region in a Venn diagram.
$A_1 - A_2$: elements in $A_1$ only, not in $A_2$. Let $x_1$ be this element. So $|A_1 - A_2| = 1$.
$A_2 - A_1$: elements in $A_2$ only, not in $A_1$. Let $x_2$ be this element. So $|A_2 - A_1| = 1$.

$A_1 - A_3$: elements in $A_1$ only, not in $A_3$. Let $y_1, y_2$ be these elements. So $|A_1 - A_3| = 2$.
$A_3 - A_1$: elements in $A_3$ only, not in $A_1$. Let $y_3, y_4$ be these elements. So $|A_3 - A_1| = 2$.

$A_2 - A_3$: elements in $A_2$ only, not in $A_3$. Let $z_1$ be this element. So $|A_2 - A_3| = 1$.
$A_3 - A_2$: elements in $A_3$ only, not in $A_2$. Let $z_2$ be this element. So $|A_3 - A_2| = 1$.

Let $X_{12} = A_1 - A_2$, $X_{21} = A_2 - A_1$.
Let $X_{13} = A_1 - A_3$, $X_{31} = A_3 - A_1$.
Let $X_{23} = A_2 - A_3$, $X_{32} = A_3 - A_2$.

Now, let's try to construct the elements systematically.
Let $a \in A_1 \cap A_2 \cap A_3$. (Common to all three)
Let $x \in (A_1 \cap A_2) - A_3$. (Common to $A_1, A_2$ but not $A_3$)
Let $y \in (A_1 \cap A_3) - A_2$. (Common to $A_1, A_3$ but not $A_2$)
Let $z \in (A_2 \cap A_3) - A_1$. (Common to $A_2, A_3$ but not $A_1$)
Let $u \in A_1 - (A_2 \cup A_3)$. (In $A_1$ only)
Let $v \in A_2 - (A_1 \cup A_3)$. (In $A_2$ only)
Let $w \in A_3 - (A_1 \cup A_2)$. (In $A_3$ only)

\begin{itemize}
    \item $A_1 - A_2 = u \cup y \cup (A_1 \cap A_3 \cap \overline{A_2}) = \{u, y\}$ must have size 1. This means these regions are not disjoint or my partitioning is wrong.

Let's use the definition of set difference directly.
$A_1 = (A_1 \setminus A_2) \cup (A_1 \cap A_2)$.
$A_2 = (A_2 \setminus A_1) \cup (A_1 \cap A_2)$.

$A_1 - A_2 = A_1 \setminus A_2 = \{x_1\}$, so $A_1$ has an element not in $A_2$. $A_2$ has an element not in $A_1$.
$A_2 - A_1 = A_2 \setminus A_1 = \{x_2\}$.

From $|A_1 - A_2| = 1$ and $|A_2 - A_1| = 1$.
This means $(A_1 \setminus A_2)$ has one element, call it $1$. So $1 \in A_1$, $1 \notin A_2$.
$(A_2 \setminus A_1)$ has one element, call it $2$. So $2 \in A_2$, $2 \notin A_1$.

From $|A_1 - A_3| = 2$ and $|A_3 - A_1| = 2$.
$A_1 \setminus A_3$ has two elements. Call them $3, 4$. So $3, 4 \in A_1$, $3, 4 \notin A_3$.
$A_3 \setminus A_1$ has two elements. Call them $5, 6$. So $5, 6 \in A_3$, $5, 6 \notin A_1$.

From $|A_2 - A_3| = 1$ and $|A_3 - A_2| = 1$.
$A_2 \setminus A_3$ has one element. Call it $7$. So $7 \in A_2$, $7 \notin A_3$.
$A_3 \setminus A_2$ has one element. Call it $8$. So $8 \in A_3$, $8 \notin A_2$.

Now, let's assemble the sets with these elements, being careful about overlaps.
$A_1$ must contain $1, 3, 4$. $A_1$ must not contain $2, 5, 6$. $A_1$ must not contain $7, 8$.
$A_2$ must contain $2, 7$. $A_2$ must not contain $1, 3, 4$. $A_2$ must not contain $5, 6, 8$.
$A_3$ must contain $5, 6, 8$. $A_3$ must not contain $1, 3, 4$. $A_3$ must not contain $2, 7$.

Let's list elements required for each portion of the Venn diagram for three sets.
Region $A_1 \setminus (A_2 \cup A_3)$: Elements only in $A_1$.
Region $A_2 \setminus (A_1 \cup A_3)$: Elements only in $A_2$.
Region $A_3 \setminus (A_1 \cup A_2)$: Elements only in $A_3$.
Region $(A_1 \cap A_2) \setminus A_3$: Elements in $A_1, A_2$ but not $A_3$.
Region $(A_1 \cap A_3) \setminus A_2$: Elements in $A_1, A_3$ but not $A_2$.
Region $(A_2 \cap A_3) \setminus A_1$: Elements in $A_2, A_3$ but not $A_1$.
Region $A_1 \cap A_2 \cap A_3$: Elements in all three.

From $A_1 \setminus A_2$: This is elements in $A_1$ but not $A_2$. This region is $A_1 \setminus (A_2 \cup A_3) \cup ((A_1 \cap A_3) \setminus A_2)$.
$|A_1 \setminus A_2| = |(A_1 \setminus (A_2 \cup A_3)) \cup ((A_1 \cap A_3) \setminus A_2)| = 1$.
$|A_2 \setminus A_1| = |(A_2 \setminus (A_1 \cup A_3)) \cup ((A_2 \cap A_3) \setminus A_1)| = 1$.
$|A_1 \setminus A_3| = |(A_1 \setminus (A_2 \cup A_3)) \cup ((A_1 \cap A_2) \setminus A_3)| = 2$.
$|A_3 \setminus A_1| = |(A_3 \setminus (A_1 \cup A_2)) \cup ((A_2 \cap A_3) \setminus A_1)| = 2$.
$|A_2 \setminus A_3| = |(A_2 \setminus (A_1 \cup A_3)) \cup ((A_1 \cap A_2) \setminus A_3)| = 1$.
$|A_3 \setminus A_2| = |(A_3 \setminus (A_1 \cup A_2)) \cup ((A_1 \cap A_3) \setminus A_2)| = 1$.

Let $p_1 = |A_1 \setminus (A_2 \cup A_3)|$, $p_2 = |A_2 \setminus (A_1 \cup A_3)|$, $p_3 = |A_3 \setminus (A_1 \cup A_2)|$.
Let $p_{12} = |(A_1 \cap A_2) \setminus A_3|$, $p_{13} = |(A_1 \cap A_3) \setminus A_2|$, $p_{23} = |(A_2 \cap A_3) \setminus A_1|$.
Let $p_{123} = |A_1 \cap A_2 \cap A_3|$.

The system of equations becomes:
\begin{align*} p_1 + p_{13} &= 1 \quad &(1) \\ p_2 + p_{23} &= 1 \quad &(2) \\ p_1 + p_{12} &= 2 \quad &(3) \\ p_3 + p_{23} &= 2 \quad &(4) \\ p_2 + p_{12} &= 1 \quad &(5) \\ p_3 + p_{13} &= 1 \quad &(6)\end{align*}
We are looking for non-negative integer solutions. Assume $p_{123}$ can be $0$ for now.

From (1) and (6): $p_1+p_{13}=1$ and $p_3+p_{13}=1$. So $p_1 = p_3$.
From (2) and (5): $p_2+p_{23}=1$ and $p_2+p_{12}=1$. So $p_{23} = p_{12}$.

Substitute $p_1=p_3$ into (4): $p_1 + p_{23} = 2$.
Substitute $p_{23}=p_{12}$ into (3): $p_1 + p_{12} = 2$. (This is consistent).

Now we have:
$p_1 + p_{13} = 1$
$p_2 + p_{12} = 1$
$p_1 + p_{12} = 2$

From $p_1 + p_{12} = 2$ and $p_2 + p_{12} = 1$:
If $p_{12} = 0$, then $p_1=2$ and $p_2=1$.
If $p_{12} = 1$, then $p_1=1$ and $p_2=0$.
If $p_{12} = 2$, then $p_1=0$ and $p_2=-1$ (not allowed).

Let's try $p_{12}=1$. Then $p_1=1$ and $p_2=0$.
Since $p_{23} = p_{12}$, $p_{23}=1$.
Since $p_1 = p_3$, $p_3=1$.
Since $p_1+p_{13}=1$, $1+p_{13}=1$, so $p_{13}=0$.

So, we have a solution for the sizes of the regions:
$p_1 = 1$ (elements only in $A_1$)
$p_2 = 0$ (elements only in $A_2$)
$p_3 = 1$ (elements only in $A_3$)
$p_{12} = 1$ (elements in $A_1, A_2$ but not $A_3$)
$p_{13} = 0$ (elements in $A_1, A_3$ but not $A_2$)
$p_{23} = 1$ (elements in $A_2, A_3$ but not $A_1$)
We can also assume $p_{123} = 0$ (elements in all three).

Let's assign specific elements based on these region sizes:
$A_1 \setminus (A_2 \cup A_3) = \{1\}$
$A_2 \setminus (A_1 \cup A_3) = \emptyset$
$A_3 \setminus (A_1 \cup A_2) = \{2\}$
$(A_1 \cap A_2) \setminus A_3 = \{3\}$
$(A_1 \cap A_3) \setminus A_2 = \emptyset$
$(A_2 \cap A_3) \setminus A_1 = \{4\}$
$A_1 \cap A_2 \cap A_3 = \emptyset$

Now construct the sets $A_1, A_2, A_3$:
$A_1 = \{1\} \cup \{3\} = \{1, 3\}$
$A_2 = \{3\} \cup \{4\} = \{3, 4\}$
$A_3 = \{2\} \cup \{4\} = \{2, 4\}$

Let's verify the difference conditions:
$A_1 - A_2 = \{1, 3\} - \{3, 4\} = \{1\}$. So $|A_1 - A_2| = 1$. (Correct)
$A_2 - A_1 = \{3, 4\} - \{1, 3\} = \{4\}$. So $|A_2 - A_1| = 1$. (Correct)
$A_1 - A_3 = \{1, 3\} - \{2, 4\} = \{1, 3\}$. So $|A_1 - A_3| = 2$. (Correct)
$A_3 - A_1 = \{2, 4\} - \{1, 3\} = \{2, 4\}$. So $|A_3 - A_1| = 2$. (Correct)
$A_2 - A_3 = \{3, 4\} - \{2, 4\} = \{3\}$. So $|A_2 - A_3| = 1$. (Correct)
$A_3 - A_2 = \{2, 4\} - \{3, 4\} = \{2\}$. So $|A_3 - A_2| = 1$. (Correct)

All conditions are met. The elements used are $\{1, 2, 3, 4\}$.
So, $A = \{1, 2, 3, 4\}$, which means $k=4$.
The sets are $A_1 = \{1, 3\}$, $A_2 = \{3, 4\}$, $A_3 = \{2, 4\}$.
The smallest $k \in \mathbb{N}$ containing these subsets is $k=4$.

\subsubsection*{Problem 1.89}
(a) For $A = \{-3, -2, \dots, 4\}$ and $B = \{1, 2, \dots, 6\}$, determine $S = \{(a, b) \in A \times B : a^2 + b^2 = 25\}$.
(b) For $C = \{a \in B : (a, b) \in S\}$ and $D = \{b \in A : (a, b) \in S\}$, where $A, B, S$ are the sets in (a), determine $C \times D$.

\paragraph*{Solution to 1.89 (a)}
$A = \{-3, -2, -1, 0, 1, 2, 3, 4\}$.
$B = \{1, 2, 3, 4, 5, 6\}$.
We need to find pairs $(a, b) \in A \times B$ such that $a^2 + b^2 = 25$.
Let's list possible $(a,b)$ pairs:
\begin{itemize}
    \item If $a = -3$, $a^2 = 9$. Then $9 + b^2 = 25 \implies b^2 = 16 \implies b = \pm 4$.
        Since $b \in B$, $b=4$ is a possible value. So $(-3, 4)$.
    \item If $a = -2$, $a^2 = 4$. Then $4 + b^2 = 25 \implies b^2 = 21$. No integer $b$.
    \item If $a = -1$, $a^2 = 1$. Then $1 + b^2 = 25 \implies b^2 = 24$. No integer $b$.
    \item If $a = 0$, $a^2 = 0$. Then $0 + b^2 = 25 \implies b^2 = 25 \implies b = \pm 5$.
        Since $b \in B$, $b=5$ is a possible value. So $(0, 5)$.
    \item If $a = 1$, $a^2 = 1$. Then $1 + b^2 = 25 \implies b^2 = 24$. No integer $b$.
    \item If $a = 2$, $a^2 = 4$. Then $4 + b^2 = 25 \implies b^2 = 21$. No integer $b$.
    \item If $a = 3$, $a^2 = 9$. Then $9 + b^2 = 25 \implies b^2 = 16 \implies b = \pm 4$.
        Since $b \in B$, $b=4$ is a possible value. So $(3, 4)$.
    \item If $a = 4$, $a^2 = 16$. Then $16 + b^2 = 25 \implies b^2 = 9 \implies b = \pm 3$.
        Since $b \in B$, $b=3$ is a possible value. So $(4, 3)$.
\end{itemize}
So, $S = \{(-3, 4), (0, 5), (3, 4), (4, 3)\}$.

\paragraph*{Solution to 1.89 (b)}
$C = \{a \in B : (a, b) \in S\}$.
This means $C$ consists of the first coordinates of pairs in $S$, but these first coordinates must be from set $B$.
The first coordinates in $S$ are $\{-3, 0, 3, 4\}$.
Elements from this set that are also in $B = \{1, 2, 3, 4, 5, 6\}$ are $\{3, 4\}$.
So, $C = \{3, 4\}$.

$D = \{b \in A : (a, b) \in S\}$.
This means $D$ consists of the second coordinates of pairs in $S$, but these second coordinates must be from set $A$.
The second coordinates in $S$ are $\{4, 5, 3\}$.
Elements from this set that are also in $A = \{-3, -2, -1, 0, 1, 2, 3, 4\}$ are $\{-3 \text{ is not a 2nd coord, } 0 \text{ is not 2nd coord}\}$. The second coordinates are $4, 5, 3$.
The elements from $\{3, 4, 5\}$ that are in $A$ are $\{3, 4\}$. Note that $5 \notin A$.
So, $D = \{3, 4\}$.

Now, determine $C \times D$:
$C \times D = \{ (3,3), (3,4), (4,3), (4,4) \}$.

\subsubsection*{Problem 1.90}
For $A = \{1, 2, 3\}$, let $B$ be the set of 2-element sets belonging to $P(A)$ and let $C$ be the set consisting of the sets that are the intersections of two distinct elements of $B$. Determine $D = P(C)$.

\paragraph*{Solution to 1.90}
First, find $P(A)$.
$A = \{1, 2, 3\}$.
$P(A) = \{\emptyset, \{1\}, \{2\}, \{3\}, \{1, 2\}, \{1, 3\}, \{2, 3\}, \{1, 2, 3\}\}$.

Next, find $B$. $B$ is the set of 2-element sets belonging to $P(A)$.
These are the subsets of $A$ that have 2 elements:
$B = \{\{1, 2\}, \{1, 3\}, \{2, 3\}\}$.

Next, find $C$. $C$ consists of the sets that are the intersections of two distinct elements of $B$.
The distinct pairs of elements from $B$ are:
\begin{itemize}
    \item $\{1, 2\} \cap \{1, 3\} = \{1\}$
    \item $\{1, 2\} \cap \{2, 3\} = \{2\}$
    \item $\{1, 3\} \cap \{2, 3\} = \{3\}$
\end{itemize}
So, $C = \{\{1\}, \{2\}, \{3\}\}$.

Finally, determine $D = P(C)$.
$|C| = 3$. So $|P(C)| = 2^3 = 8$.
$P(C) = \{\emptyset, \{\{1\}\}, \{\{2\}\}, \{\{3\}\}, \{\{1\}, \{2\}\}, \{\{1\}, \{3\}\}, \{\{2\}, \{3\}\}, \{\{1\}, \{2\}, \{3\}\}\}$.
So, $D = P(C)$.

\subsubsection*{Problem 1.91}
For a real number $r$, let $A_r = \{r, r+1\}$. Let $S = \{x \in \mathbb{R} : x^2 + 2x - 1 = 0\}$.
(a) Determine $B = A_s \times A_t$, for the distinct elements $s, t \in S$, where $s < t$.
(b) Let $C = \{ab : (a, b) \in B\}$. Determine the sum of the elements of $C$.

\paragraph*{Solution to 1.91}
First, find the elements of $S$.
$S = \{x \in \mathbb{R} : x^2 + 2x - 1 = 0\}$.
Using the quadratic formula $x = \frac{-b \pm \sqrt{b^2 - 4ac}}{2a}$:
$x = \frac{-2 \pm \sqrt{2^2 - 4(1)(-1)}}{2(1)} = \frac{-2 \pm \sqrt{4 + 4}}{2} = \frac{-2 \pm \sqrt{8}}{2} = \frac{-2 \pm 2\sqrt{2}}{2} = -1 \pm \sqrt{2}$.
So $S = \{-1 - \sqrt{2}, -1 + \sqrt{2}\}$.
The distinct elements are $s = -1 - \sqrt{2}$ and $t = -1 + \sqrt{2}$, with $s < t$.

\paragraph*{Solution to 1.91 (a)}
Determine $B = A_s \times A_t$.
$A_s = \{s, s+1\} = \{-1 - \sqrt{2}, (-1 - \sqrt{2}) + 1\} = \{-1 - \sqrt{2}, -\sqrt{2}\}$.
$A_t = \{t, t+1\} = \{-1 + \sqrt{2}, (-1 + \sqrt{2}) + 1\} = \{-1 + \sqrt{2}, \sqrt{2}\}$.

$B = A_s \times A_t = \{(-1 - \sqrt{2}, -1 + \sqrt{2}), (-1 - \sqrt{2}, \sqrt{2}), (-\sqrt{2}, -1 + \sqrt{2}), (-\sqrt{2}, \sqrt{2})\}$.

\paragraph*{Solution to 1.91 (b)}
Let $C = \{ab : (a, b) \in B\}$. Determine the sum of the elements of $C$.
The elements of $C$ are the products of the components of each ordered pair in $B$.
1. $(-1 - \sqrt{2})(-1 + \sqrt{2}) = ((-1)^2 - (\sqrt{2})^2) = 1 - 2 = -1$.
2. $(-1 - \sqrt{2})(\sqrt{2}) = -\sqrt{2} - (\sqrt{2})^2 = -\sqrt{2} - 2$.
3. $(-\sqrt{2})(-1 + \sqrt{2}) = \sqrt{2} - (\sqrt{2})^2 = \sqrt{2} - 2$.
4. $(-\sqrt{2})(\sqrt{2}) = -(\sqrt{2})^2 = -2$.

So $C = \{-1, -2 - \sqrt{2}, -2 + \sqrt{2}, -2\}$.
Note that $-2$ appears twice. In a set, elements are distinct. So $C = \{-1, -2 - \sqrt{2}, -2 + \sqrt{2}, -2\}$.
The sum of the elements in $C$:
Sum $= (-1) + (-2 - \sqrt{2}) + (-2 + \sqrt{2}) + (-2)$
Sum $= -1 - 2 - \sqrt{2} - 2 + \sqrt{2} - 2$
Sum $= -1 - 2 - 2 - 2 = -7$.

\subsubsection*{Problem 1.92}
(a) For a set $A$ with $|A| = 2$, what is the largest possible value of $|A \cap P(A)|$?
(b) What is the largest possible value of $|A \cap P(A)|$ if $|A| = 3$?
(c) This should suggest another question to you. What is the answer to your question?

\paragraph*{Solution to 1.92 (a)}
Let $A$ be a set with $|A| = 2$. Let $A = \{x, y\}$.
Then $P(A) = \{\emptyset, \{x\}, \{y\}, \{x, y\}\}$.
We want to find the largest possible value of $|A \cap P(A)|$.
The elements of $A$ are $x$ and $y$.
The elements of $P(A)$ are $\emptyset, \{x\}, \{y\}, \{x, y\}$.
For an element of $A$ to be in $P(A)$, it must be a subset of $A$.
\begin{itemize}
    \item Is $x \in P(A)$? This means $x$ must be one of $\emptyset, \{x\}, \{y\}, \{x, y\}$. This happens if $x$ is equal to $\{x\}$, or $x$ is equal to $\{y\}$, or $x$ is equal to $\emptyset$, or $x$ is equal to $\{x,y\}$. This is generally not true unless $x$ is itself a set (and one of these specific sets). For example, if $x=\{\text{apple}\}$, then $x \ne \{x\}$.
    \item Is $x = \emptyset$? Not for typical elements.
    \item Is $x = \{x\}$? This is a non-standard set theory concept (a set containing itself as an element), usually disallowed by axioms of ZFC.
    \item Is $x = \{y\}$? This implies $x$ is a set and contains only $y$.
    \item Is $x = \{x, y\}$? This implies $x$ is a set and contains $x$ and $y$.
\end{itemize}
However, the problem implicitly assumes we can construct $A$ in such a way. For $A \cap P(A)$ to be non-empty, some elements of $A$ must *also* be subsets of $A$.
Let $A = \{x, \{x\}\}$. Here $|A|=2$.
$P(A) = \{\emptyset, \{x\}, \{\{x\}\}, \{x, \{x\}\}\}$.
Then $A \cap P(A) = \{\{x\}\}$. So $|A \cap P(A)| = 1$.
What if $A = \{\emptyset, \{x\}\}$? No, this makes $|A|$ not 2.
What if $A = \{\emptyset, \{ \}\}$. This is the same as $A = \{\emptyset\}$. This makes $|A|=1$.

Let $A = \{e_1, e_2\}$. We want $e_1 \in P(A)$ or $e_2 \in P(A)$.
If $e_1 = \emptyset$, then $A = \{\emptyset, e_2\}$. Then $\emptyset \in P(A)$. So $A \cap P(A)$ contains $\emptyset$.
Let $A = \{\emptyset, \{1\}\}$. Here $|A|=2$.
$P(A) = \{\emptyset, \{\emptyset\}, \{\{1\}\}, \{\emptyset, \{1\}\}\}$.
$A \cap P(A) = \{\emptyset\}$. So $|A \cap P(A)| = 1$.

What if $e_2 = e_1$? Then $|A|=1$, not 2.
Could $e_1 = \{e_2\}$ and $e_2 = \{e_1\}$? This is a cyclic definition.
Could $e_1 = \{e_1\}$? Usually avoided.

The maximum possible value of $|A \cap P(A)|$ for $|A|=n$ is related to how many elements of $A$ can be subsets of $A$.
For $|A|=2$, $A=\{e_1, e_2\}$.
Subsets of $A$ are $\emptyset, \{e_1\}, \{e_2\}, \{e_1, e_2\}$.
We want $e_1$ to be one of these, or $e_2$ to be one of these.
If $e_1 = \emptyset$, then $A = \{\emptyset, e_2\}$. $\emptyset \in A \cap P(A)$.
If $e_2 = \{e_1\}$, then $A = \{e_1, \{e_1\}\}$.
If $e_1 = \emptyset$ AND $e_2 = \{e_1\} = \{\emptyset\}$.
Then $A = \{\emptyset, \{\emptyset\}\}$. This satisfies $|A|=2$.
$P(A) = \{\emptyset, \{\emptyset\}, \{\{\emptyset\}\}, \{\emptyset, \{\emptyset\}\}\}$.
$A \cap P(A) = \{\emptyset, \{\emptyset\}\}$.
In this case, $|A \cap P(A)| = 2$. This is the largest possible value, as we cannot have more elements in the intersection than in $A$ itself.

\paragraph*{Solution to 1.92 (b)}
If $|A| = 3$. Let $A = \{e_1, e_2, e_3\}$.
We want to maximize the number of elements of $A$ that are also subsets of $A$.
Let $A = \{\emptyset, \{\emptyset\}, \{\emptyset, \{\emptyset\}\}\}$.
$|A|=3$.
$P(A)$ contains:
$\emptyset$
$\{\emptyset\}$
$\{\{\emptyset\}\}$
$\{\emptyset, \{\emptyset\}\}$
$\{\emptyset, \{\emptyset\}\}$ (this is $A_2$, the second element)
$\{\emptyset, \{\emptyset, \{\emptyset\}\}\}$
$\{\{\emptyset\}, \{\emptyset, \{\emptyset\}\}\}$
$A$ itself.
Elements of $A$:
\begin{enumerate}
    \item $\emptyset$: $\emptyset \in P(A)$.
    \item $\{\emptyset\}$: $\{\emptyset\} \in P(A)$.
    \item $\{\emptyset, \{\emptyset\}\}$: This element is $A$ itself. $A \in P(A)$.
\end{enumerate}
So, all three elements of $A$ are also subsets of $A$.
Thus, for $A = \{\emptyset, \{\emptyset\}, \{\emptyset, \{\emptyset\}\}\}$, we have $A \cap P(A) = A$.
So, $|A \cap P(A)| = 3$.
The largest possible value is 3.

\paragraph*{Solution to 1.92 (c)}
This suggests the question: For a set $A$ with $|A|=n$, what is the largest possible value of $|A \cap P(A)|$?
The answer is $n$.
This can be achieved by constructing $A$ in an iterative way (a specific type of construction based on the empty set).
Let $A_0 = \emptyset$.
Let $A_1 = \{\emptyset\}$.
Let $A_2 = \{\emptyset, \{\emptyset\}\}$.
Let $A_3 = \{\emptyset, \{\emptyset\}, \{\emptyset, \{\emptyset\}\}\}$.
In general, for $|A|=n$, we can construct $A$ such that $A = \{A_0, A_1, \dots, A_{n-1}\}$, where $A_k = P(A_{k-1})$ in a specific way such that each $A_i$ is built upon the previous and becomes an element.
More simply, $A_0 = \emptyset$, $A_{k+1} = A_k \cup \{A_k\}$. This is not quite correct for all elements to be subsets.
A more standard way to define sets like this is using von Neumann ordinals, where each ordinal is the set of all smaller ordinals.
$0 = \emptyset$
$1 = \{\emptyset\}$
$2 = \{\emptyset, \{\emptyset\}\}$
$3 = \{\emptyset, \{\emptyset\}, \{\emptyset, \{\emptyset\}\}\}$
If $A$ is chosen as the von Neumann ordinal $n$, then $A$ contains $n$ elements, and all these elements are also subsets of $A$.
For example, if $A = 3 = \{\emptyset, \{\emptyset\}, \{\emptyset, \{\emptyset\}\}\}$.
Its elements are $\emptyset$, $\{\emptyset\}$, and $\{\emptyset, \{\emptyset\}\}$.
$\emptyset$ is a subset of any set, so $\emptyset \in P(A)$.
$\{\emptyset\}$ is a subset of $A$ because $\emptyset \in A$. So $\{\emptyset\} \in P(A)$.
$\{\emptyset, \{\emptyset\}\}$ is a subset of $A$ because $\emptyset \in A$ and $\{\emptyset\} \in A$. So $\{\emptyset, \{\emptyset\}\} \in P(A)$.
Thus, all elements of $A$ are also subsets of $A$.
So, $|A \cap P(A)| = |A| = n$.
The largest possible value of $|A \cap P(A)|$ for $|A|=n$ is $n$.

\subsubsection*{Problem 1.93}
For $n \in \mathbb{N}$, let $S_n = \{x + y : x,y \in \mathbb{R}, \frac{1}{n^2} \le x^2 + y^2 \le 1\}$. Determine $\bigcup_{n=1}^\infty S_n$ and $\bigcap_{n=1}^\infty S_n$.

\paragraph*{Solution to 1.93}
The set $S_n$ is defined by sums of $x$ and $y$ where $(x,y)$ lies in an annulus (ring) in the $\mathbb{R}^2$ plane.
Specifically, $x^2 + y^2$ is between $\frac{1}{n^2}$ and $1$. This means the points $(x,y)$ are between (and including) circles of radius $\frac{1}{n}$ and radius $1$, centered at the origin.

Let $K_n = \{(x,y) \in \mathbb{R}^2 : \frac{1}{n^2} \le x^2 + y^2 \le 1\}$.
The range of $x$ and $y$ values depends on the annulus.
For any $(x,y) \in K_n$, we want to find the range of $x+y$.
Let $x+y = k$. We are looking for the set of all possible $k$.
Consider $x^2+y^2=R^2$. The line $x+y=k$ intersects this circle.
$y = k-x$.
$x^2 + (k-x)^2 = R^2$
$x^2 + k^2 - 2kx + x^2 = R^2$
$2x^2 - 2kx + k^2 - R^2 = 0$.
For real solutions for $x$, the discriminant must be non-negative:
$(-2k)^2 - 4(2)(k^2 - R^2) \ge 0$
$4k^2 - 8k^2 + 8R^2 \ge 0$
$8R^2 - 4k^2 \ge 0$
$8R^2 \ge 4k^2 \implies 2R^2 \ge k^2 \implies |k| \le \sqrt{2}R$.
So for a fixed radius $R$, $x+y$ can range from $-\sqrt{2}R$ to $\sqrt{2}R$.

In $K_n$, $R^2$ ranges from $\frac{1}{n^2}$ to $1$. So $R$ ranges from $\frac{1}{n}$ to $1$.
Therefore, $x+y$ can range from $-\sqrt{2} \cdot 1$ to $\sqrt{2} \cdot 1$, and also from $-\sqrt{2} \cdot \frac{1}{n}$ to $\sqrt{2} \cdot \frac{1}{n}$.
The set $S_n$ is the interval $[-\sqrt{2} \cdot 1, \sqrt{2} \cdot 1] = [-\sqrt{2}, \sqrt{2}]$.
This is incorrect. The range of $x+y$ for points in $K_n$ is determined by the minimum and maximum possible values of $x+y$.
The maximum value of $x+y$ occurs when $x=y$ and $x^2+y^2$ is maximized.
If $x^2+y^2=1$, then $x=y=\frac{1}{\sqrt{2}}$. Then $x+y = \sqrt{2}$.
If $x^2+y^2=\frac{1}{n^2}$, then $x=y=\frac{1}{n\sqrt{2}}$. Then $x+y = \frac{\sqrt{2}}{n}$.
So $x+y$ can range from $-\sqrt{2}$ to $\sqrt{2}$.
And $x+y$ can range from $-\frac{\sqrt{2}}{n}$ to $\frac{\sqrt{2}}{n}$.
The set $S_n = [-\sqrt{2}, \sqrt{2}]$ for all $n \in \mathbb{N}$, because the range of $x^2+y^2$ is $[1/n^2, 1]$.
No, this is not true. $S_n$ is the union of all ranges $[-\sqrt{2}R, \sqrt{2}R]$ for $R \in [\frac{1}{n}, 1]$.
As $R$ increases from $\frac{1}{n}$ to $1$, the interval $[-\sqrt{2}R, \sqrt{2}R]$ expands.
The smallest interval is for $R = \frac{1}{n}$, giving $[-\frac{\sqrt{2}}{n}, \frac{\sqrt{2}}{n}]$.
The largest interval is for $R=1$, giving $[-\sqrt{2}, \sqrt{2}]$.
So, $S_n = [-\sqrt{2}, \sqrt{2}]$.

Let's re-verify the reasoning that $S_n = [-\sqrt{2}, \sqrt{2}]$.
Let $k = x+y$. We found that for $(x,y)$ on a circle of radius $R$, $k \in [-\sqrt{2}R, \sqrt{2}R]$.
The set $K_n$ is the union of all circles with radius $R$ where $R \in [\frac{1}{n}, 1]$.
So $S_n = \bigcup_{R \in [\frac{1}{n}, 1]} [-\sqrt{2}R, \sqrt{2}R]$.
The union of these increasing intervals is the largest interval, which occurs at $R=1$.
So $S_n = [-\sqrt{2}, \sqrt{2}]$.

Now for the union and intersection:
$\bigcup_{n=1}^\infty S_n = \bigcup_{n=1}^\infty [-\sqrt{2}, \sqrt{2}] = [-\sqrt{2}, \sqrt{2}]$.
$\bigcap_{n=1}^\infty S_n = \bigcap_{n=1}^\infty [-\sqrt{2}, \sqrt{2}] = [-\sqrt{2}, \sqrt{2}]$.

This seems too simple, let me check the definition of $S_n$ again.
Is it possible that $x+y$ cannot take all values in $[-\sqrt{2}R, \sqrt{2}R]$ for points in $K_n$?
The $S_n$ description is "set of all $x+y$".
Let $f(x,y) = x+y$. We are looking for the image of $K_n$ under $f$.
The set $K_n$ is a connected region in $\mathbb{R}^2$. The function $f(x,y)$ is continuous.
Thus, the image $f(K_n)$ must be a connected interval.
The minimum value of $x+y$ over $K_n$ is $-\sqrt{2}R_{max} = -\sqrt{2}$.
The maximum value of $x+y$ over $K_n$ is $\sqrt{2}R_{max} = \sqrt{2}$.
So $S_n = [-\sqrt{2}, \sqrt{2}]$ for all $n \in \mathbb{N}$.

Therefore, the union and intersection are indeed:
$\bigcup_{n=1}^\infty S_n = [-\sqrt{2}, \sqrt{2}]$.
$\bigcap_{n=1}^\infty S_n = [-\sqrt{2}, \sqrt{2}]$.

\subsubsection*{Problem 1.94}
For the set $S = \{1, 2, \dots, 7\}$, consider the five subsets
$A_1 = \{1, 2\}$, $A_2 = \{3, 4, 5\}$, $A_3 = \{2, 5, 6\}$, $A_4 = \{2, 7\}$, $A_5 = \{1, 2\}$.

\paragraph*{Solution to 1.94 (a)}
Is the collection $P_1 = \{A_1, A_2, A_3, A_4, A_5\}$ a partition of $S$?
Recall the conditions for a partition:
1.  All subsets must be nonempty. (All $A_i$ are nonempty.)
2.  The union of all subsets must be $S$.
3.  The subsets must be pairwise disjoint (except for equal sets, which count as one).

Let's check condition 3 first.
$A_1 = \{1, 2\}$, $A_5 = \{1, 2\}$. So $A_1 = A_5$. This means the collection effectively has 4 distinct sets: $\{A_1, A_2, A_3, A_4\}$.
Are they pairwise disjoint?
$A_1 \cap A_2 = \{1, 2\} \cap \{3, 4, 5\} = \emptyset$.
$A_1 \cap A_3 = \{1, 2\} \cap \{2, 5, 6\} = \{2\} \neq \emptyset$.
Since $A_1 \cap A_3 \neq \emptyset$, the sets are not pairwise disjoint.
Therefore, $P_1$ is not a partition of $S$.

\paragraph*{Solution to 1.94 (b)}
Let $B_i = A_i - A_{i+1}$ for $i = 1, 2, 3, 4$. Is $P_2 = \{B_1, B_2, B_3, B_4\}$ a partition of $S$?
First, let's calculate each $B_i$:
$B_1 = A_1 - A_2 = \{1, 2\} - \{3, 4, 5\} = \{1, 2\}$.
$B_2 = A_2 - A_3 = \{3, 4, 5\} - \{2, 5, 6\} = \{3, 4\}$.
$B_3 = A_3 - A_4 = \{2, 5, 6\} - \{2, 7\} = \{5, 6\}$.
$B_4 = A_4 - A_5 = \{2, 7\} - \{1, 2\} = \{7\}$.

So $P_2 = \{\{1, 2\}, \{3, 4\}, \{5, 6\}, \{7\}\}$.
Let's check the partition conditions for $P_2$:
1.  All subsets are nonempty. (Yes, each $B_i$ contains elements.)
2.  The subsets are pairwise disjoint.
    \begin{itemize}
        \item $\{1, 2\} \cap \{3, 4\} = \emptyset$
        \item $\{1, 2\} \cap \{5, 6\} = \emptyset$
        \item $\{1, 2\} \cap \{7\} = \emptyset$
        \item $\{3, 4\} \cap \{5, 6\} = \emptyset$
        \item $\{3, 4\} \cap \{7\} = \emptyset$
        \item $\{5, 6\} \cap \{7\} = \emptyset$
    \end{itemize}
    (Yes, they are pairwise disjoint.)
3.  The union of all subsets must be $S = \{1, 2, 3, 4, 5, 6, 7\}$.
    $B_1 \cup B_2 \cup B_3 \cup B_4 = \{1, 2\} \cup \{3, 4\} \cup \{5, 6\} \cup \{7\} = \{1, 2, 3, 4, 5, 6, 7\}$.
    (Yes, the union is $S$.)
All three conditions are met. Therefore, $P_2$ is a partition of $S$.

\paragraph*{Solution to 1.94 (c)}
Let $C_i = A_{6-i} - A_{5-i}$ for $i = 1, 2, 3, 4$. Is $P_3 = \{C_1, C_2, C_3, C_4\}$ a partition of $S$?
Let's calculate each $C_i$:
$C_1 = A_{6-1} - A_{5-1} = A_5 - A_4 = \{1, 2\} - \{2, 7\} = \{1\}$.
$C_2 = A_{6-2} - A_{5-2} = A_4 - A_3 = \{2, 7\} - \{2, 5, 6\} = \{7\}$.
$C_3 = A_{6-3} - A_{5-3} = A_3 - A_2 = \{2, 5, 6\} - \{3, 4, 5\} = \{2, 6\}$.
$C_4 = A_{6-4} - A_{5-4} = A_2 - A_1 = \{3, 4, 5\} - \{1, 2\} = \{3, 4, 5\}$.

So $P_3 = \{\{1\}, \{7\}, \{2, 6\}, \{3, 4, 5\}\}$.
Let's check the partition conditions for $P_3$:
1.  All subsets are nonempty. (Yes.)
2.  The subsets are pairwise disjoint.
    \begin{itemize}
        \item $\{1\} \cap \{7\} = \emptyset$
        \item $\{1\} \cap \{2, 6\} = \emptyset$
        \item $\{1\} \cap \{3, 4, 5\} = \emptyset$
        \item $\{7\} \cap \{2, 6\} = \emptyset$
        \item $\{7\} \cap \{3, 4, 5\} = \emptyset$
        \item $\{2, 6\} \cap \{3, 4, 5\} = \emptyset$
    \end{itemize}
    (Yes, they are pairwise disjoint.)
3.  The union of all subsets must be $S = \{1, 2, 3, 4, 5, 6, 7\}$.
    $C_1 \cup C_2 \cup C_3 \cup C_4 = \{1\} \cup \{7\} \cup \{2, 6\} \cup \{3, 4, 5\} = \{1, 2, 3, 4, 5, 6, 7\}$.
    (Yes, the union is $S$.)
All three conditions are met. Therefore, $P_3$ is a partition of $S$.

\subsubsection*{Problem 1.95}
Let $S = \{1, 2, 3\}$. Find a condition under which nonempty subsets $A, B, C$ of $S$ have the property that $|A \times (B \cup C)| = |A \times B| + |A \times C|$.

\paragraph*{Solution to 1.95}
We know that for any finite sets $X, Y, Z$:
$|X \times Y| = |X| \cdot |Y|$.
So the given property is $|A| \cdot |B \cup C| = |A| \cdot |B| + |A| \cdot |C|$.
Since $A$ is a nonempty subset of $S$, $|A| \neq 0$. We can divide by $|A|$:
$|B \cup C| = |B| + |C|$.

We also know the Principle of Inclusion-Exclusion for two sets:
$|B \cup C| = |B| + |C| - |B \cap C|$.
Comparing this with the desired condition, $|B \cup C| = |B| + |C|$, we must have $|B \cap C| = 0$.
The cardinal number of $B \cap C$ is 0 if and only if $B \cap C = \emptyset$.
Therefore, the condition is that $B$ and $C$ must be disjoint.
The problem also states that $A, B, C$ are nonempty subsets of $S$.

So, the condition is that $B$ and $C$ are nonempty and disjoint.
Since $B$ and $C$ are subsets of $S$, and are nonempty and disjoint.
Example: Let $S=\{1,2,3\}$. $A=\{1\}$. $B=\{2\}$. $C=\{3\}$.
$B \cup C = \{2,3\}$. $|B \cup C| = 2$.
$|A \times (B \cup C)| = |\{1\} \times \{2,3\}| = |\{(1,2), (1,3)\}| = 2$.
$|A \times B| = |\{1\} \times \{2\}| = |\{(1,2)\}| = 1$.
$|A \times C| = |\{1\} \times \{3\}| = |\{(1,3)\}| = 1$.
$|A \times B| + |A \times C| = 1 + 1 = 2$.
The equality holds. This example has $B \cap C = \emptyset$.

\section{In-Depth Analysis of Problems and Techniques}

\subsection{A) Problem Types and General Approach}
The homework problems cover a range of fundamental set theory concepts, requiring careful application of definitions.

\begin{enumerate}
    \item \textbf{Describing Sets (Listing/Set-Builder):}
        *   \textbf{Problems 1.73, 1.74, 1.75:} These problems test the ability to translate between two primary methods of set description: listing elements and using set-builder notation.
        *   \textbf{General Approach:} For listing, substitute values (especially integers) into the defining expression or solve the given equation/inequality for integer or natural number domains. For set-builder notation, identify the common property or range of elements from the given list. Pay close attention to the domain of the variable (e.g., $\mathbb{Z}, \mathbb{N}, \mathbb{R}, \mathbb{Q}$).
    \item \textbf{Cardinality Calculation:}
        *   \textbf{Problem 1.76:} Focuses on determining the number of distinct elements in a set. This includes sets with nested elements (sets as elements themselves) and empty sets.
        *   \textbf{General Approach:} Carefully identify each unique element. Remember that repeated listings of the same object count as one element. For sets defined by conditions, solve the condition first to find the elements, then count them. For special sets like the empty set, the cardinality is 0.
    \item \textbf{Cartesian Products:}
        *   \textbf{Problems 1.77, 1.78(d), 1.89(a), 1.91(a), 1.95:} These problems require constructing the set of all ordered pairs from two given sets.
        *   \textbf{General Approach:} Systematically pair each element from the first set with every element from the second set, forming $(a,b)$ tuples. Remember that $(a,b) \neq (b,a)$ unless $a=b$.
    \item \textbf{Set Operations (Union, Intersection, Difference, Complement):}
        *   \textbf{Problems 1.78(a, b, c), 1.83, 1.94(b, c), 1.95:} These are core problems involving the basic operations.
        *   \textbf{General Approach:} Apply the definitions step-by-step. For expressions with multiple operations, follow order of operations (often parentheses first). For complements, identify the universal set. For problems involving conditions, apply logical "and", "or", "not" to the conditions. Use Venn diagrams for visualization if helpful.
    \item \textbf{Power Sets and Subsets:}
        *   \textbf{Problems 1.79, 1.80, 1.81, 1.90, 1.92:} These problems involve constructing power sets and understanding the relationship between sets and their subsets.
        *   \textbf{General Approach:} List all possible combinations of elements, including the empty set and the set itself. For $P(A)$, each element is a subset of $A$. Be careful when elements of a set are themselves sets (nested sets); they are treated as single, distinct objects for the purpose of forming subsets. Distinguish between an element $\in$ and a subset $\subseteq$.
    \item \textbf{Equality of Sets:}
        *   \textbf{Problem 1.82:} Tests the understanding that two sets are equal if and only if they contain exactly the same elements.
        *   \textbf{General Approach:} List all elements for each set, then compare the lists.
    \item \textbf{Working with Indexed Collections of Sets:}
        *   \textbf{Problems 1.84, 1.93:} These problems involve finding the union or intersection of an infinite collection of sets, often defined by a parameter.
        *   \textbf{General Approach:} Analyze the pattern of how the sets $A_r$ (or $S_n$) change as the index $r$ (or $n$) varies. For unions, identify the largest possible range covered by any set in the collection. For intersections, find the common elements across *all* sets in the collection. For continuous parameters, consider limits as the parameter approaches its bounds.
    \item \textbf{Partition Problems:}
        *   \textbf{Problem 1.94(a, b, c):} Requires checking if a collection of subsets satisfies the three conditions for a partition.
        *   \textbf{General Approach:} Verify that all subsets are nonempty, they are pairwise disjoint, and their union forms the original set.
    \item \textbf{Constructive/Proof-based Problems:}
        *   \textbf{Problems 1.79, 1.80, 1.81, 1.85, 1.86, 1.88, 1.92, 1.95:} These problems ask to create examples of sets or prove conditions based on set properties.
        *   \textbf{General Approach:} Start with simple sets and build up. For conditions involving cardinalities, ensure distinct elements are assigned to different regions of a Venn diagram as needed. For proofs, use element-wise arguments (e.g., "let $x$ be an element of set $X$, then...").

\subsection{B) Key Algebraic and Calculus Manipulations}
\begin{enumerate}
    \item \textbf{Solving Equations and Inequalities:} Essential for determining elements of sets defined by conditions.
        *   \textbf{Examples:}
            *   $x^3 - 4x = 0$ (Problem 1.75(a), 1.82): Factoring polynomials ($x(x^2-4)=0$) to find integer or real roots.
            *   $x^2 + 2x - 1 = 0$ (Problem 1.91): Using the quadratic formula to find real roots ($x = -1 \pm \sqrt{2}$).
            *   $1 \le k^2 < 100$ (Problem 1.75(g), 1.76(e,f)): Solving inequalities involving squares, considering both positive and negative roots.
            *   Absolute value inequalities like $|x|=-1$ (Problem 1.75(b), 1.76(b)) are critical to identify empty sets.
    \item \textbf{Interpretation of Domains ($\mathbb{N}, \mathbb{Z}, \mathbb{Q}, \mathbb{R}$):} Understanding which numbers are allowed is crucial.
        *   \textbf{Example:} $2x+2=0 \implies x=-1$. If $x \in \mathbb{N}$ (Problem 1.82 B), the set is $\emptyset$. If $x \in \mathbb{Z}$, it's $\{-1\}$.
    \item \textbf{Distinguishing Elements vs. Subsets:} This is a conceptual and notational hurdle that often trips up students, especially with power sets.
        *   \textbf{Example:} In $A = \{1, 2, \{1, 2\}\}$, $1 \in A$ but $\{1\} \notin A$. However, $\{1\}$ is a subset of $A$. For $P(A)$, $\{1\} \in P(A)$.
        *   \textbf{Problem 1.92:} Explicitly addresses this distinction by asking how many elements of a set can also be subsets of that set.
    \item \textbf{Careful Processing of Set-Builder Notation and Nested Sets:}
        *   \textbf{Example:} $A = \{1, 2, 3, \{1, 2, 3\}, 4, \{4\}\}$ (Problem 1.76(a)). The element $\{1, 2, 3\}$ is a single element, distinct from $1, 2, 3$.
    \item \textbf{Systematic Approach for Set Operations:} Performing operations like union, intersection, and difference step-by-step for complex expressions.
        *   \textbf{Example:} $(A \cup B) - (B \cap C)$ (Problem 1.78(a)). Calculate $A \cup B$ first, then $B \cap C$, then the difference.
    \item \textbf{Decomposition for Venn Diagram Type Problems:} For problems like 1.88, breaking down set differences into disjoint regions of a Venn diagram and setting up a system of equations for the cardinalities of these regions helps in constructing the sets. This often involves linear algebra (solving systems of equations).
\end{enumerate}

\section{"Cheatsheet" and Tips for Success}

\subsection{Most Important Formulas}
\begin{itemize}
    \item \textbf{Cardinality of a finite set $A$:} $|A|$ (number of distinct elements)
    \item \textbf{Cardinality of the Power Set of $A$:} $|P(A)| = 2^{|A|}$
    \item \textbf{Union:} $A \cup B = \{x \mid x \in A \text{ or } x \in B\}$
    \item \textbf{Intersection:} $A \cap B = \{x \mid x \in A \text{ and } x \in B\}$
    \item \textbf{Set Difference:} $A - B = \{x \mid x \in A \text{ and } x \notin B\}$
    \item \textbf{Complement (given universal set $U$):} $\bar{A} = U - A = \{x \mid x \in U \text{ and } x \notin A\}$
    \item \textbf{Cartesian Product:} $A \times B = \{(a, b) \mid a \in A \text{ and } b \in B\}$
    \item \textbf{Cardinality of Cartesian Product:} $|A \times B| = |A| \cdot |B|$
    \item \textbf{Principle of Inclusion-Exclusion (for 2 sets):} $|A \cup B| = |A| + |B| - |A \cap B|$
    \item \textbf{Partition Properties:} A collection of nonempty subsets of $A$ forms a partition if they are pairwise disjoint and their union is $A$.
\end{itemize}

\subsection{Cheats, Tricks, and Shortcuts}
\begin{itemize}
    \item \textbf{The Empty Set is Universal:} $\emptyset$ is a subset of every set. Its cardinality is 0.
    \item \textbf{Order Doesn't Matter (for sets):} $\{1, 2, 3\}$ is the same as $\{3, 1, 2\}$.
    \item \textbf{Order DOES Matter (for ordered pairs):} $(1, 2)$ is NOT the same as $(2, 1)$ unless $1=2$.
    \item \textbf{Venn Diagrams:} Use them to visualize set operations, especially for 2 or 3 sets. They can help clarify disjoint regions and unions/intersections.
    \item \textbf{Element-by-Element Analysis:} When proving set equalities (e.g., $A \subseteq B$), pick an arbitrary element $x \in A$ and show it must be in $B$.
    \item \textbf{Simplify Conditions:} Before listing elements, simplify any algebraic expressions or inequalities within set-builder notation.
    \item \textbf{Domain Check:} Always double-check the domain of elements (e.g., $x \in \mathbb{Z}$, $x \in \mathbb{R}$). This frequently determines if a set is empty or finite.
\end{itemize}

\subsection{Common Pitfalls, Mistakes, or Traps to Watch Out For}
\begin{itemize}
    \item \textbf{Confusing $\in$ (is an element of) with $\subseteq$ (is a subset of):} This is the most common error.
        *   Example: For $A = \{1, \{2\}\}$, $1 \in A$ and $\{2\} \in A$. But $\{1\} \subseteq A$ and $\{\{2\}\} \subseteq A$. $1 \not\subseteq A$.
    \item \textbf{Misinterpreting the Empty Set:}
        *   $\emptyset$ is a set with no elements.
        *   $\{\emptyset\}$ is a set with one element (which happens to be the empty set). Its cardinality is 1.
    \item \textbf{Cardinality of Nested Sets:} Each distinct object, regardless of whether it's a number or another set, counts as one element for cardinality.
        *   Example: $|\{1, \{2, 3\}, 4\}| = 3$. The element $\{2, 3\}$ is one element.
    \item \textbf{Universal Set for Complements:} Always be clear about what the universal set $U$ is when finding a complement.
    \item \textbf{Contradictory Problem Statements:} As seen in Problem 1.83, sometimes conditions provided can be inconsistent, leading to no possible solution.
    \item \textbf{Order in Cartesian Products:} Remember that $(a,b)$ is an ordered pair.
\end{itemize}

\subsection{Tips for How to Recognize Which Type of Problem You Are Dealing With}
\begin{itemize}
    \item \textbf{Keywords:} "List elements," "describe using set-builder notation" (describing sets). "Cardinality," "number of elements" (cardinality). "Subset," "proper subset," "power set" (subsets/power sets). "Union," "intersection," "difference," "complement," "disjoint" (set operations). "Cartesian product," "ordered pairs" (Cartesian products). "Partition" (partitions).
    \item \textbf{Notation:} Curly braces for sets, $\in, \notin, \subseteq, \subset, \cup, \cap, -, \bar{A}, \times, |A|$.
    \item \textbf{Implicit Structures:} Problems asking for constructing sets with specific properties often involve working backward or using a systematic trial-and-error approach combined with Venn diagram logic.
    \item \textbf{Infinite Collections:} Problems with $\bigcup_{i=1}^\infty$ or $\bigcap_{i \in I}$ indicate indexed families of sets, requiring analysis of patterns or limits.
\end{itemize}

\section{Conceptual Synthesis and The "Big Picture"}

\subsection{A) Thematic Connections}
The core theme of set theory is the rigorous definition and manipulation of \textbf{collections of objects}. This fundamental idea provides the bedrock for almost all other branches of mathematics. Sets allow us to precisely define mathematical structures and relationships.

\begin{itemize}
    \item \textbf{Number Systems:} Integers, rationals, reals, and complex numbers are themselves defined as sets with specific operations and properties. For example, the natural numbers can be constructed purely from set theory (e.g., von Neumann ordinals where $0 = \emptyset$, $1 = \{0\}$, $2 = \{0, 1\}$, etc.).
    \item \textbf{Functions and Relations:} A function is formally defined as a special type of relation, which in turn is a set of ordered pairs (a subset of a Cartesian product). This connection allows for a precise understanding of mapping between sets.
    \item \textbf{Logic:} The "and," "or," and "not" operations in logic have direct counterparts in set theory as intersection, union, and complement, respectively. This highlights the deep connection between propositional logic and set theory.
    \item \textbf{Probability Theory:} A sample space in probability is defined as the set of all possible outcomes of an experiment. Events are subsets of this sample space. Set operations (union, intersection, complement) are directly used to describe compound events.
\item \textbf{Topology and Analysis:} Concepts like open sets, closed sets, and neighborhoods are defined using sets of real numbers or points in metric spaces, which are fundamentally set-theoretic constructs.
\end{itemize}

\subsection{B) Forward and Backward Links}
\begin{itemize}
    \item \textbf{Backward Links (Prerequisites):}
        *   **Basic Arithmetic and Algebra:** Essential for solving equations and inequalities that define set elements.
        *   **Propositional Logic:** The foundations of set operations (union, intersection, complement) are directly analogous to logical "or," "and," and "not."
        *   **Number Systems:** Understanding $\mathbb{N}, \mathbb{Z}, \mathbb{Q}, \mathbb{R}, \mathbb{C}$ is crucial because sets often draw their elements from these domains.
    \item \textbf{Forward Links (Future Topics):}
        *   **Relations and Functions:** These are defined as specific types of sets of ordered pairs (subsets of Cartesian products). Understanding Cartesian products and set-builder notation is vital.
        *   **Probability and Statistics:** The concepts of sample spaces, events, mutually exclusive events, and unions/intersections of events are all deeply rooted in set theory. Set theory provides the language to define and analyze randomness.
        *   **Abstract Algebra:** Groups, rings, fields, and vector spaces are defined as sets equipped with specific binary operations that satisfy certain axioms. Set theory provides the very definition of these fundamental algebraic structures.
        *   **Topology and Real Analysis:** These advanced fields heavily rely on precise definitions of sets (open sets, closed sets, compact sets, connected sets) to study properties of spaces and functions.
        *   **Computer Science:** Data structures like lists, arrays, and hash sets are practical implementations of sets. Set operations are fundamental in database queries and algorithm design.
        *   **Mathematical Finance/Economics/Stochastic Calculus:**
            *   **Event Spaces:** In stochastic calculus and derivative pricing, set theory defines the underlying sample space of all possible future states of the world or asset prices. Events (e.g., "stock price above a certain level," "interest rates rise") are subsets of this sample space.
            *   **Portfolio Theory:** A portfolio can be viewed as a set of financial assets. Operations like combining portfolios (union), finding common assets (intersection), or diversifying by removing assets (set difference) leverage set-theoretic concepts.
            *   **Utility Theory:** Consumer preferences or agent choices are often modeled over sets of goods or consumption bundles.
            *   **Option Payoffs:** The conditions for an option's payoff (e.g., call option payoff is $\max(S_T - K, 0)$) implicitly define sets of outcomes ($S_T > K$ vs $S_T \le K$) where different functions apply.

\section{Real-World Application and Modeling}

\subsection{A) Concrete Scenarios (Finance/Stats/Econ focus)}
1.  \textbf{Portfolio Diversification and Overlap Analysis:} A fund manager has several investment portfolios. Each portfolio is a set of assets (stocks, bonds, derivatives). Set operations are essential for understanding the composition, overlap, and diversification of these portfolios. For example, the union of two portfolios shows all assets held across both, while the intersection reveals common holdings. The difference shows assets unique to one portfolio. This helps in managing risk and identifying redundant or missing investments.
2.  \textbf{Market Segmentation and Customer Targeting in Finance:} A retail bank wants to target specific customer groups for new financial products (e.g., high-yield savings accounts, personal loans). They segment their customer base into sets based on demographics, income, credit score, and product history. Using set intersections (e.g., customers with high income AND high credit score) allows for precise targeting, while unions help broaden reach.
3.  \textbf{Event Modeling in Quantitative Finance (Stochastic Processes):} In modeling financial markets, possible future scenarios (e.g., stock price movements, interest rate changes) form event spaces. Set theory is used to define complex events. For example, an event "stock A rises AND stock B falls" is an intersection of two simpler events (sets of market states). This is foundational for Monte Carlo simulations and risk analysis where probabilities are assigned to these sets.

\subsection{B) Model Problem Setup}
\textbf{Scenario: Portfolio Diversification and Risk Management}

A boutique investment firm manages two equity portfolios, Portfolio Alpha and Portfolio Beta, for its clients. The firm is considering merging these portfolios or analyzing their unique holdings to understand overall exposure.

\textbf{Mathematical Model:}
Let $U$ be the universal set of all publicly traded stocks in a specific market (e.g., NYSE and NASDAQ).
Let $P_\alpha$ be the set of stocks held in Portfolio Alpha.
Let $P_\beta$ be the set of stocks held in Portfolio Beta.

\begin{itemize}
    \item \textbf{Problem Description:} The firm wants to identify:
        \begin{enumerate}
            \item All unique stocks held across both portfolios.
            \item Stocks common to both portfolios (potential overlap/redundancy).
            \item Stocks held only in Portfolio Alpha (unique to Alpha).
            \item Stocks held only in Portfolio Beta (unique to Beta).
            \item Stocks NOT held in either portfolio that might be candidates for future investment (complement relative to $U$).
        \end{enumerate}
    \item \textbf{Variables:}
        *   $U$: A large, finite set of all investable stocks, e.g., $U = \{\text{AAPL, GOOG, MSFT, TSLA, JPM, GS, ...}\}$.
        *   $P_\alpha$: A finite subset of $U$ (stocks in Alpha).
        *   $P_\beta$: A finite subset of $U$ (stocks in Beta).
    \item \textbf{Formulation (Set Operations):}
        \begin{enumerate}
            \item All unique stocks across both portfolios: $P_\alpha \cup P_\beta$. The number of such stocks is $|P_\alpha \cup P_\beta|$.
            \item Stocks common to both portfolios: $P_\alpha \cap P_\beta$. The number of such stocks is $|P_\alpha \cap P_\beta|$.
            \item Stocks held only in Portfolio Alpha: $P_\alpha - P_\beta$. The number of such stocks is $|P_\alpha - P_\beta|$.
            \item Stocks held only in Portfolio Beta: $P_\beta - P_\alpha$. The number of such stocks is $|P_\beta - P_\alpha|$.
            \item Stocks not held in either portfolio: $\overline{P_\alpha \cup P_\beta} = U - (P_\alpha \cup P_\beta)$. The number of such stocks is $|U - (P_\alpha \cup P_\beta)|$.
        \end{enumerate}
\end{itemize}
This model uses basic set operations to analyze and manage investment portfolios, helping to inform decisions on diversification, risk, and asset allocation.

\section{Common Variations and Untested Concepts}

The homework problems provide a solid foundation in basic set definitions and operations. However, a comprehensive understanding of set theory extends to several key concepts not directly covered.

\subsection{1. De Morgan's Laws}
\subsubsection*{Explanation}
De Morgan's Laws are fundamental identities that relate the complement of a union or intersection of sets to the union or intersection of their complements. They are crucial for simplifying set expressions and are direct analogues of De Morgan's Laws in propositional logic.
\begin{itemize}
    \item $\overline{A \cup B} = \bar{A} \cap \bar{B}$ (The complement of the union is the intersection of the complements).
    \item $\overline{A \cap B} = \bar{A} \cup \bar{B}$ (The complement of the intersection is the union of the complements).
\end{itemize}

\subsubsection*{Worked-Out Example}
Let $U = \{1, 2, 3, 4, 5\}$, $A = \{1, 2, 3\}$, $B = \{3, 4\}$.
We will verify the first law: $\overline{A \cup B} = \bar{A} \cap \bar{B}$.

\paragraph*{Left Hand Side (LHS): $\overline{A \cup B}$}
First, find $A \cup B$:
$A \cup B = \{1, 2, 3\} \cup \{3, 4\} = \{1, 2, 3, 4\}$.
Next, find its complement relative to $U$:
$\overline{A \cup B} = U - (A \cup B) = \{1, 2, 3, 4, 5\} - \{1, 2, 3, 4\} = \{5\}$.

\paragraph*{Right Hand Side (RHS): $\bar{A} \cap \bar{B}$}
First, find $\bar{A}$:
$\bar{A} = U - A = \{1, 2, 3, 4, 5\} - \{1, 2, 3\} = \{4, 5\}$.
Next, find $\bar{B}$:
$\bar{B} = U - B = \{1, 2, 3, 4, 5\} - \{3, 4\} = \{1, 2, 5\}$.
Finally, find their intersection:
$\bar{A} \cap \bar{B} = \{4, 5\} \cap \{1, 2, 5\} = \{5\}$.

Since LHS = RHS ($ \{5\} = \{5\} $), De Morgan's first law is verified for this example.

\subsection{2. Principle of Inclusion-Exclusion (for Cardinality of Union)}
\subsubsection*{Explanation}
While the homework uses the idea implicitly, the formal principle provides a way to calculate the cardinality of a union of sets by adding individual cardinalities, subtracting overlaps, adding triple overlaps, and so on. For two sets, it's $|A \cup B| = |A| + |B| - |A \cap B|$. For three sets, it expands to:
$|A \cup B \cup C| = |A| + |B| + |C| - (|A \cap B| + |A \cap C| + |B \cap C|) + |A \cap B \cap C|$.
This principle is invaluable in combinatorics and probability.

\subsubsection*{Worked-Out Example}
Let $A = \{1, 2, 3, 4\}$, $B = \{3, 4, 5, 6\}$, and $C = \{2, 3, 5, 7\}$.
We want to find $|A \cup B \cup C|$ using the Principle of Inclusion-Exclusion.

First, find individual cardinalities:
$|A| = 4$
$|B| = 4$
$|C| = 4$

Next, find cardinalities of pairwise intersections:
$A \cap B = \{3, 4\}$, so $|A \cap B| = 2$.
$A \cap C = \{2, 3\}$, so $|A \cap C| = 2$.
$B \cap C = \{3, 5\}$, so $|B \cap C| = 2$.

Finally, find cardinality of the triple intersection:
$A \cap B \cap C = \{1, 2, 3, 4\} \cap \{3, 4, 5, 6\} \cap \{2, 3, 5, 7\} = \{3\}$.
So $|A \cap B \cap C| = 1$.

Now, apply the formula:
$|A \cup B \cup C| = |A| + |B| + |C| - (|A \cap B| + |A \cap C| + |B \cap C|) + |A \cap B \cap C|$
$|A \cup B \cup C| = 4 + 4 + 4 - (2 + 2 + 2) + 1$
$|A \cup B \cup C| = 12 - 6 + 1 = 7$.

\paragraph*{Verification by direct listing:}
$A \cup B \cup C = \{1, 2, 3, 4\} \cup \{3, 4, 5, 6\} \cup \{2, 3, 5, 7\} = \{1, 2, 3, 4, 5, 6, 7\}$.
Indeed, $|A \cup B \cup C| = 7$.

\subsection{3. Russell's Paradox}
\subsubsection*{Explanation}
Russell's Paradox, discovered by Bertrand Russell in 1901, revealed a fundamental flaw in "naive" set theory (where any well-defined collection can form a set). It demonstrated that assuming such a principle leads to contradictions. The paradox arises from considering "the set of all sets that do not contain themselves."

\subsubsection*{Worked-Out Example}
Let $R$ be the set of all sets that are not members of themselves.
$R = \{X \mid X \text{ is a set and } X \notin X\}$.

Now, consider the question: Is $R$ a member of itself ($R \in R$)?
\begin{itemize}
    \item \textbf{Case 1: Assume $R \in R$.}
        If $R \in R$, then by the defining property of $R$, $R$ must be a set that does not contain itself. So, $R \notin R$.
        This is a contradiction ($R \in R$ and $R \notin R$).
    \item \textbf{Case 2: Assume $R \notin R$.}
        If $R \notin R$, then $R$ satisfies the defining property of the set $R$ (it is a set that does not contain itself). Therefore, $R$ must be a member of $R$. So, $R \in R$.
        This is also a contradiction ($R \notin R$ and $R \in R$).
\end{itemize}
Since both assumptions lead to a contradiction, the very definition of such a set $R$ leads to a paradox. This necessitated the development of axiomatic set theories (like ZFC) that restrict which collections can form sets, thus avoiding such self-referential paradoxes.

\section{Advanced Diagnostic Testing: "Find the Flaw"}

\subsection{Problem 1}
**Problem:** Let $A = \{1, \{1\}\}$. Determine $|P(A)|$.
**Flawed Solution:**
Given $A = \{1, \{1\}\}$.
The elements of $A$ are $1$ and $1$.
So $|A|=1$.
Therefore, $|P(A)| = 2^{|A|} = 2^1 = 2$.
$P(A) = \{\emptyset, \{1\}\}$.

\subsection{Problem 2}
**Problem:** Let $U = \{1, 2, 3, 4, 5\}$, $A = \{1, 3\}$, $B = \{3, 5\}$. Determine $A \cup \bar{B}$.
**Flawed Solution:**
First, find $\bar{B}$. $\bar{B} = \{1, 2, 4\}$.
Then, $A \cup \bar{B} = \{1, 3\} \cup \{1, 2, 4\} = \{1, 2, 3, 4, 1\}$.
So $A \cup \bar{B} = \{1, 2, 3, 4\}$.

\subsection{Problem 3}
**Problem:** Let $A = \{x \in \mathbb{Z} : x^2 < 5\}$. List the elements of $A$.
**Flawed Solution:**
$x^2 < 5$. Taking the square root of both sides gives $x < \sqrt{5}$.
Since $\sqrt{5} \approx 2.23$, the integers less than $2.23$ are $2, 1, 0, -1, -2, \dots$.
So $A = \{\dots, -2, -1, 0, 1, 2\}$.

\subsection{Problem 4}
**Problem:** For sets $X = \{a, b\}$ and $Y = \{c, d\}$, determine $Y \times X$.
**Flawed Solution:**
$Y \times X = \{(a,c), (a,d), (b,c), (b,d)\}$.

\subsection{Problem 5}
**Problem:** Let $S = \{1, 2, 3, 4\}$. Is $P = \{\{1, 2\}, \{2, 3\}, \{4\}\}$ a partition of $S$?
**Flawed Solution:**
1. All subsets are nonempty. (True)
2. The union of all subsets is $S$. $\{1, 2\} \cup \{2, 3\} \cup \{4\} = \{1, 2, 3, 4\} = S$. (True)
Since conditions 1 and 2 are met, $P$ is a partition of $S$.

\end{document}2 are met, $P$ is a partition of $S$.

\end{document}